\chapter{Bizonyítás különböző modellekben}
\label{chap:proofs}

\begin{code}[hide]%
\>[0]\AgdaSymbol{\{-\#}\AgdaSpace{}%
\AgdaKeyword{OPTIONS}\AgdaSpace{}%
\AgdaPragma{--prop}\AgdaSpace{}%
\AgdaSymbol{\#-\}}\<%
\\
%
\\[\AgdaEmptyExtraSkip]%
\>[0]\AgdaKeyword{open}\AgdaSpace{}%
\AgdaKeyword{import}\AgdaSpace{}%
\AgdaModule{Agda.Primitive}\<%
\\
\>[0]\AgdaKeyword{open}\AgdaSpace{}%
\AgdaKeyword{import}\AgdaSpace{}%
\AgdaModule{Agda.Builtin.Bool}\<%
\\
\>[0]\AgdaKeyword{open}\AgdaSpace{}%
\AgdaKeyword{import}\AgdaSpace{}%
\AgdaModule{Lib}\<%
\\
%
\\[\AgdaEmptyExtraSkip]%
\>[0]\AgdaKeyword{data}\AgdaSpace{}%
\AgdaDatatype{funar}\AgdaSpace{}%
\AgdaSymbol{:}\AgdaSpace{}%
\AgdaDatatype{ℕ}\AgdaSpace{}%
\AgdaSymbol{→}\AgdaSpace{}%
\AgdaPrimitive{Set}\AgdaSpace{}%
\AgdaKeyword{where}\<%
\\
\>[0][@{}l@{\AgdaIndent{0}}]%
\>[2]\AgdaInductiveConstructor{zero}\AgdaSpace{}%
\AgdaSymbol{:}\AgdaSpace{}%
\AgdaDatatype{funar}\AgdaSpace{}%
\AgdaNumber{0}\<%
\\
%
\>[2]\AgdaInductiveConstructor{suc}%
\>[7]\AgdaSymbol{:}\AgdaSpace{}%
\AgdaDatatype{funar}\AgdaSpace{}%
\AgdaNumber{1}\<%
\\
%
\>[2]\AgdaInductiveConstructor{+'}%
\>[7]\AgdaSymbol{:}\AgdaSpace{}%
\AgdaDatatype{funar}\AgdaSpace{}%
\AgdaNumber{2}\<%
\\
%
\\[\AgdaEmptyExtraSkip]%
\>[0]\AgdaKeyword{data}\AgdaSpace{}%
\AgdaDatatype{relar}\AgdaSpace{}%
\AgdaSymbol{:}\AgdaSpace{}%
\AgdaDatatype{ℕ}\AgdaSpace{}%
\AgdaSymbol{→}\AgdaSpace{}%
\AgdaPrimitive{Set}\AgdaSpace{}%
\AgdaKeyword{where}\<%
\\
\>[0][@{}l@{\AgdaIndent{0}}]%
\>[2]\AgdaInductiveConstructor{eq}%
\>[7]\AgdaSymbol{:}\AgdaSpace{}%
\AgdaDatatype{relar}\AgdaSpace{}%
\AgdaNumber{2}\<%
\\
%
\\[\AgdaEmptyExtraSkip]%
\>[0]\AgdaKeyword{open}\AgdaSpace{}%
\AgdaKeyword{import}\AgdaSpace{}%
\AgdaModule{model}\<%
\\
%
\\[\AgdaEmptyExtraSkip]%
\>[0]\AgdaKeyword{module}\AgdaSpace{}%
\AgdaModule{workInAnyModel}\AgdaSpace{}%
\AgdaSymbol{\{}\AgdaBound{i}\AgdaSymbol{\}\{}\AgdaBound{j}\AgdaSymbol{\}\{}\AgdaBound{k}\AgdaSymbol{\}\{}\AgdaBound{l}\AgdaSymbol{\}(}\AgdaBound{M}\AgdaSpace{}%
\AgdaSymbol{:}\AgdaSpace{}%
\AgdaRecord{Model}\AgdaSpace{}%
\AgdaDatatype{funar}\AgdaSpace{}%
\AgdaDatatype{relar}\AgdaSpace{}%
\AgdaSymbol{\{}\AgdaBound{i}\AgdaSymbol{\}\{}\AgdaBound{j}\AgdaSymbol{\}\{}\AgdaBound{k}\AgdaSymbol{\}\{}\AgdaBound{l}\AgdaSymbol{\})}\AgdaSpace{}%
\AgdaKeyword{where}\<%
\\
\>[0][@{}l@{\AgdaIndent{0}}]%
\>[2]\AgdaKeyword{open}\AgdaSpace{}%
\AgdaModule{Model}\AgdaSpace{}%
\AgdaBound{M}\<%
\\
\>[0]\<%
\\
%
\>[2]\AgdaFunction{zero'}\AgdaSpace{}%
\AgdaSymbol{:}\AgdaSpace{}%
\AgdaSymbol{∀\{}\AgdaBound{Γ}\AgdaSymbol{\}}\AgdaSpace{}%
\AgdaSymbol{→}\AgdaSpace{}%
\AgdaField{Tm}\AgdaSpace{}%
\AgdaBound{Γ}\<%
\\
%
\>[2]\AgdaFunction{zero'}\AgdaSpace{}%
\AgdaSymbol{\{}\AgdaBound{Γ}\AgdaSymbol{\}}\AgdaSpace{}%
\AgdaSymbol{=}\AgdaSpace{}%
\AgdaField{fun}\AgdaSpace{}%
\AgdaInductiveConstructor{zero}\AgdaSpace{}%
\AgdaSymbol{\AgdaUnderscore{}}\<%
\\
%
\\[\AgdaEmptyExtraSkip]%
%
\>[2]\AgdaFunction{suc'}\AgdaSpace{}%
\AgdaSymbol{:}\AgdaSpace{}%
\AgdaSymbol{∀\{}\AgdaBound{Γ}\AgdaSymbol{\}}\AgdaSpace{}%
\AgdaSymbol{→}\AgdaSpace{}%
\AgdaField{Tm}\AgdaSpace{}%
\AgdaBound{Γ}\AgdaSpace{}%
\AgdaSymbol{→}\AgdaSpace{}%
\AgdaField{Tm}\AgdaSpace{}%
\AgdaBound{Γ}\<%
\\
%
\>[2]\AgdaFunction{suc'}\AgdaSpace{}%
\AgdaSymbol{\{}\AgdaBound{Γ}\AgdaSymbol{\}}\AgdaSpace{}%
\AgdaBound{n}\AgdaSpace{}%
\AgdaSymbol{=}\AgdaSpace{}%
\AgdaField{fun}\AgdaSpace{}%
\AgdaSymbol{\{}\AgdaBound{Γ}\AgdaSymbol{\}}\AgdaSpace{}%
\AgdaInductiveConstructor{suc}\AgdaSpace{}%
\AgdaSymbol{(}\AgdaBound{n}\AgdaSpace{}%
\AgdaOperator{\AgdaInductiveConstructor{,}}\AgdaSpace{}%
\AgdaSymbol{\AgdaUnderscore{})}\<%
\\
%
\\[\AgdaEmptyExtraSkip]%
%
\>[2]\AgdaFunction{three}\AgdaSpace{}%
\AgdaSymbol{:}\AgdaSpace{}%
\AgdaSymbol{∀\{}\AgdaBound{Γ}\AgdaSymbol{\}}\AgdaSpace{}%
\AgdaSymbol{→}\AgdaSpace{}%
\AgdaField{Tm}\AgdaSpace{}%
\AgdaBound{Γ}\<%
\\
%
\>[2]\AgdaFunction{three}\AgdaSpace{}%
\AgdaSymbol{\{}\AgdaBound{Γ}\AgdaSymbol{\}}\AgdaSpace{}%
\AgdaSymbol{=}\AgdaSpace{}%
\AgdaFunction{suc'}\AgdaSpace{}%
\AgdaSymbol{\{}\AgdaBound{Γ}\AgdaSymbol{\}}\AgdaSpace{}%
\AgdaSymbol{(}\AgdaFunction{suc'}\AgdaSpace{}%
\AgdaSymbol{\{}\AgdaBound{Γ}\AgdaSymbol{\}}\AgdaSpace{}%
\AgdaSymbol{(}\AgdaFunction{suc'}\AgdaSpace{}%
\AgdaSymbol{\{}\AgdaBound{Γ}\AgdaSymbol{\}}\AgdaSpace{}%
\AgdaSymbol{(}\AgdaFunction{zero'}\AgdaSpace{}%
\AgdaSymbol{\{}\AgdaBound{Γ}\AgdaSymbol{\})))}\<%
\\
%
\\[\AgdaEmptyExtraSkip]%
%
\>[2]\AgdaFunction{Eq}\AgdaSpace{}%
\AgdaSymbol{:}\AgdaSpace{}%
\AgdaSymbol{∀\{}\AgdaBound{Γ}\AgdaSymbol{\}}\AgdaSpace{}%
\AgdaSymbol{→}\AgdaSpace{}%
\AgdaField{Tm}\AgdaSpace{}%
\AgdaBound{Γ}\AgdaSpace{}%
\AgdaSymbol{→}\AgdaSpace{}%
\AgdaField{Tm}\AgdaSpace{}%
\AgdaBound{Γ}\AgdaSpace{}%
\AgdaSymbol{→}\AgdaSpace{}%
\AgdaField{For}\AgdaSpace{}%
\AgdaBound{Γ}\<%
\\
%
\>[2]\AgdaFunction{Eq}\AgdaSpace{}%
\AgdaSymbol{\{}\AgdaBound{Γ}\AgdaSymbol{\}}\AgdaSpace{}%
\AgdaBound{u}\AgdaSpace{}%
\AgdaBound{v}\AgdaSpace{}%
\AgdaSymbol{=}\AgdaSpace{}%
\AgdaField{Rel}\AgdaSpace{}%
\AgdaSymbol{\{}\AgdaBound{Γ}\AgdaSymbol{\}}\AgdaSpace{}%
\AgdaInductiveConstructor{eq}\AgdaSpace{}%
\AgdaSymbol{(}\AgdaBound{u}\AgdaSpace{}%
\AgdaOperator{\AgdaInductiveConstructor{,}}\AgdaSpace{}%
\AgdaSymbol{(}\AgdaBound{v}\AgdaSpace{}%
\AgdaOperator{\AgdaInductiveConstructor{,}}\AgdaSpace{}%
\AgdaSymbol{\AgdaUnderscore{}))}\<%
\\
%
\\[\AgdaEmptyExtraSkip]%
%
\>[2]\AgdaFunction{A⊃A}\AgdaSpace{}%
\AgdaSymbol{:}\AgdaSpace{}%
\AgdaSymbol{\{}\AgdaBound{A}\AgdaSpace{}%
\AgdaSymbol{:}\AgdaSpace{}%
\AgdaField{For}\AgdaSpace{}%
\AgdaField{◇}\AgdaSymbol{\}}\AgdaSpace{}%
\AgdaSymbol{→}\AgdaSpace{}%
\AgdaField{Pf}\AgdaSpace{}%
\AgdaSymbol{(}\AgdaField{◇}\AgdaSymbol{)}\AgdaSpace{}%
\AgdaSymbol{(}\AgdaBound{A}\AgdaSpace{}%
\AgdaOperator{\AgdaField{⊃}}\AgdaSpace{}%
\AgdaBound{A}\AgdaSymbol{)}\<%
\\
%
\>[2]\AgdaFunction{A⊃A}\AgdaSpace{}%
\AgdaSymbol{=}\AgdaSpace{}%
\AgdaField{⊃in}\AgdaSpace{}%
\AgdaField{qₚ}\<%
\end{code}
\newcommand{\pAnyE}{
\begin{code}%
%
\>[2]\AgdaFunction{A⊃B⊃A}\AgdaSpace{}%
\AgdaSymbol{:}\AgdaSpace{}%
\AgdaSymbol{\{}\AgdaBound{A}\AgdaSpace{}%
\AgdaBound{B}\AgdaSpace{}%
\AgdaSymbol{:}\AgdaSpace{}%
\AgdaField{For}\AgdaSpace{}%
\AgdaField{◇}\AgdaSymbol{\}}\AgdaSpace{}%
\AgdaSymbol{→}\AgdaSpace{}%
\AgdaField{Pf}\AgdaSpace{}%
\AgdaSymbol{(}\AgdaField{◇}\AgdaSymbol{)}\AgdaSpace{}%
\AgdaSymbol{(}\AgdaBound{A}\AgdaSpace{}%
\AgdaOperator{\AgdaField{⊃}}\AgdaSpace{}%
\AgdaBound{B}\AgdaSpace{}%
\AgdaOperator{\AgdaField{⊃}}\AgdaSpace{}%
\AgdaBound{A}\AgdaSymbol{)}\<%
\\
%
\>[2]\AgdaFunction{A⊃B⊃A}\AgdaSpace{}%
\AgdaSymbol{\{}\AgdaBound{A}\AgdaSymbol{\}\{}\AgdaBound{B}\AgdaSymbol{\}}\AgdaSpace{}%
\AgdaSymbol{=}\<%
\\
\>[2][@{}l@{\AgdaIndent{0}}]%
\>[4]\AgdaField{⊃in}%
\>[137I]\AgdaSymbol{(}\<%
\\
\>[.][@{}l@{}]\<[137I]%
\>[8]\AgdaFunction{substP}\<%
\\
\>[8][@{}l@{\AgdaIndent{0}}]%
\>[12]\AgdaSymbol{(λ}\AgdaSpace{}%
\AgdaBound{x}\AgdaSpace{}%
\AgdaSymbol{→}\AgdaSpace{}%
\AgdaField{Pf}\AgdaSpace{}%
\AgdaSymbol{(}\AgdaField{◇}\AgdaSpace{}%
\AgdaOperator{\AgdaField{▹ₚ}}\AgdaSpace{}%
\AgdaBound{A}\AgdaSymbol{)}\AgdaSpace{}%
\AgdaSymbol{(}\AgdaBound{x}\AgdaSymbol{))}\<%
\\
%
\>[12]\AgdaSymbol{(}\AgdaField{⊃[]}\AgdaSpace{}%
\AgdaOperator{\AgdaFunction{⁻¹}}\AgdaSymbol{)}\<%
\\
%
\>[12]\AgdaSymbol{(}\AgdaField{⊃in}\AgdaSpace{}%
\AgdaSymbol{\{}\AgdaField{◇}\AgdaSpace{}%
\AgdaOperator{\AgdaField{▹ₚ}}\AgdaSpace{}%
\AgdaBound{A}\AgdaSymbol{\}\{}\AgdaBound{B}\AgdaSpace{}%
\AgdaOperator{\AgdaField{[}}\AgdaSpace{}%
\AgdaField{pₚ}\AgdaSpace{}%
\AgdaOperator{\AgdaField{]ᶠ}}\AgdaSymbol{\}\{}\AgdaBound{A}\AgdaSpace{}%
\AgdaOperator{\AgdaField{[}}\AgdaSpace{}%
\AgdaField{pₚ}\AgdaSpace{}%
\AgdaOperator{\AgdaField{]ᶠ}}\AgdaSymbol{\}}\AgdaSpace{}%
\AgdaSymbol{(}\AgdaField{qₚ}\AgdaSpace{}%
\AgdaOperator{\AgdaField{[}}\AgdaSpace{}%
\AgdaField{pₚ}\AgdaSpace{}%
\AgdaOperator{\AgdaField{]ᵖ}}\AgdaSymbol{))}\<%
\\
%
\>[4]\AgdaSymbol{)}\<%
\end{code}}
\begin{code}[hide]%
%
\>[2]\AgdaFunction{A∧B⊃B∨A}\AgdaSpace{}%
\AgdaSymbol{:}\AgdaSpace{}%
\AgdaSymbol{\{}\AgdaBound{A}\AgdaSpace{}%
\AgdaBound{B}\AgdaSpace{}%
\AgdaSymbol{:}\AgdaSpace{}%
\AgdaField{For}\AgdaSpace{}%
\AgdaField{◇}\AgdaSymbol{\}}\AgdaSpace{}%
\AgdaSymbol{→}\AgdaSpace{}%
\AgdaField{Pf}\AgdaSpace{}%
\AgdaSymbol{(}\AgdaField{◇}\AgdaSymbol{)}\AgdaSpace{}%
\AgdaSymbol{(}\AgdaBound{A}\AgdaSpace{}%
\AgdaOperator{\AgdaField{∧}}\AgdaSpace{}%
\AgdaBound{B}\AgdaSpace{}%
\AgdaOperator{\AgdaField{⊃}}\AgdaSpace{}%
\AgdaBound{B}\AgdaSpace{}%
\AgdaOperator{\AgdaField{∨}}\AgdaSpace{}%
\AgdaBound{A}\AgdaSymbol{)}\<%
\\
%
\>[2]\AgdaFunction{A∧B⊃B∨A}\AgdaSpace{}%
\AgdaSymbol{\{}\AgdaBound{A}\AgdaSymbol{\}\{}\AgdaBound{B}\AgdaSymbol{\}}\AgdaSpace{}%
\AgdaSymbol{=}\AgdaSpace{}%
\AgdaField{⊃in}\AgdaSpace{}%
\AgdaSymbol{(}\AgdaFunction{substP}\AgdaSpace{}%
\AgdaSymbol{(λ}\AgdaSpace{}%
\AgdaBound{x}\AgdaSpace{}%
\AgdaSymbol{→}\AgdaSpace{}%
\AgdaField{Pf}\AgdaSpace{}%
\AgdaSymbol{(}\AgdaField{◇}\AgdaSpace{}%
\AgdaOperator{\AgdaField{▹ₚ}}\AgdaSpace{}%
\AgdaBound{A}\AgdaSpace{}%
\AgdaOperator{\AgdaField{∧}}\AgdaSpace{}%
\AgdaBound{B}\AgdaSymbol{)}\AgdaSpace{}%
\AgdaBound{x}\AgdaSymbol{)}\AgdaSpace{}%
\AgdaSymbol{(}\AgdaField{∨[]}\AgdaSpace{}%
\AgdaOperator{\AgdaFunction{⁻¹}}\AgdaSymbol{)}\AgdaSpace{}%
\AgdaSymbol{(}\AgdaField{∨in₁}\AgdaSpace{}%
\AgdaSymbol{(}\AgdaField{∧out₂}\AgdaSpace{}%
\AgdaSymbol{(}\AgdaFunction{substP}\AgdaSpace{}%
\AgdaSymbol{(λ}\AgdaSpace{}%
\AgdaBound{x}\AgdaSpace{}%
\AgdaSymbol{→}\AgdaSpace{}%
\AgdaField{Pf}\AgdaSpace{}%
\AgdaSymbol{(}\AgdaField{◇}\AgdaSpace{}%
\AgdaOperator{\AgdaField{▹ₚ}}\AgdaSpace{}%
\AgdaBound{A}\AgdaSpace{}%
\AgdaOperator{\AgdaField{∧}}\AgdaSpace{}%
\AgdaBound{B}\AgdaSymbol{)}\AgdaSpace{}%
\AgdaBound{x}\AgdaSymbol{)}\AgdaSpace{}%
\AgdaField{∧[]}\AgdaSpace{}%
\AgdaField{qₚ}\AgdaSymbol{))))}\<%
\end{code}
\newcommand{\pAnyC}{
\begin{code}%
%
\>[2]\AgdaFunction{∀P∧Q⊃∀P∧∀Q}\AgdaSpace{}%
\AgdaSymbol{:}%
\>[207I]\AgdaSymbol{\{}\AgdaBound{P}\AgdaSpace{}%
\AgdaBound{Q}\AgdaSpace{}%
\AgdaSymbol{:}\AgdaSpace{}%
\AgdaField{For}\AgdaSpace{}%
\AgdaSymbol{(}\AgdaField{◇}\AgdaSpace{}%
\AgdaOperator{\AgdaField{▹ₜ}}\AgdaSymbol{)\}}\AgdaSpace{}%
\AgdaSymbol{→}\<%
\\
\>[207I][@{}l@{\AgdaIndent{0}}]%
\>[16]\AgdaField{Pf}\AgdaSpace{}%
\AgdaSymbol{(}\AgdaField{◇}\AgdaSymbol{)}\AgdaSpace{}%
\AgdaSymbol{((}\AgdaField{Forall}\AgdaSpace{}%
\AgdaSymbol{(}\AgdaBound{P}\AgdaSpace{}%
\AgdaOperator{\AgdaField{∧}}\AgdaSpace{}%
\AgdaBound{Q}\AgdaSymbol{))}\AgdaSpace{}%
\AgdaOperator{\AgdaField{⊃}}\AgdaSpace{}%
\AgdaSymbol{(}\AgdaField{Forall}\AgdaSpace{}%
\AgdaBound{P}\AgdaSpace{}%
\AgdaOperator{\AgdaField{∧}}\AgdaSpace{}%
\AgdaField{Forall}\AgdaSpace{}%
\AgdaBound{Q}\AgdaSymbol{))}\<%
\\
%
\>[2]\AgdaFunction{∀P∧Q⊃∀P∧∀Q}\AgdaSpace{}%
\AgdaSymbol{\{}\AgdaBound{P}\AgdaSymbol{\}\{}\AgdaBound{Q}\AgdaSymbol{\}}\AgdaSpace{}%
\AgdaSymbol{=}\<%
\\
\>[2][@{}l@{\AgdaIndent{0}}]%
\>[4]\AgdaField{⊃in}\<%
\\
%
\>[4]\AgdaSymbol{(}\AgdaFunction{substP}\AgdaSpace{}%
\AgdaSymbol{(λ}\AgdaSpace{}%
\AgdaBound{x}\AgdaSpace{}%
\AgdaSymbol{→}\AgdaSpace{}%
\AgdaField{Pf}\AgdaSpace{}%
\AgdaSymbol{(}\AgdaField{◇}\AgdaSpace{}%
\AgdaOperator{\AgdaField{▹ₚ}}\AgdaSpace{}%
\AgdaField{Forall}\AgdaSpace{}%
\AgdaSymbol{(}\AgdaBound{P}\AgdaSpace{}%
\AgdaOperator{\AgdaField{∧}}\AgdaSpace{}%
\AgdaBound{Q}\AgdaSymbol{))}\AgdaSpace{}%
\AgdaBound{x}\AgdaSymbol{)}\AgdaSpace{}%
\AgdaSymbol{(}\AgdaField{∧[]}\AgdaSpace{}%
\AgdaOperator{\AgdaFunction{⁻¹}}\AgdaSymbol{)}\<%
\\
\>[4][@{}l@{\AgdaIndent{0}}]%
\>[6]\AgdaSymbol{(}\AgdaField{∧in}\<%
\\
\>[6][@{}l@{\AgdaIndent{0}}]%
\>[8]\AgdaSymbol{(}\AgdaFunction{substP}\AgdaSpace{}%
\AgdaSymbol{(λ}\AgdaSpace{}%
\AgdaBound{x}\AgdaSpace{}%
\AgdaSymbol{→}\AgdaSpace{}%
\AgdaField{Pf}\AgdaSpace{}%
\AgdaSymbol{(}\AgdaField{◇}\AgdaSpace{}%
\AgdaOperator{\AgdaField{▹ₚ}}\AgdaSpace{}%
\AgdaField{Forall}\AgdaSpace{}%
\AgdaSymbol{(}\AgdaBound{P}\AgdaSpace{}%
\AgdaOperator{\AgdaField{∧}}\AgdaSpace{}%
\AgdaBound{Q}\AgdaSymbol{))}\AgdaSpace{}%
\AgdaBound{x}\AgdaSymbol{)}\AgdaSpace{}%
\AgdaSymbol{(}\AgdaField{Forall[]}\AgdaSpace{}%
\AgdaOperator{\AgdaFunction{⁻¹}}\AgdaSymbol{)}\<%
\\
\>[8][@{}l@{\AgdaIndent{0}}]%
\>[10]\AgdaSymbol{(}\AgdaField{∀in}\AgdaSpace{}%
\AgdaSymbol{(}\AgdaField{∧out₁}\AgdaSpace{}%
\AgdaSymbol{(}\AgdaField{∀out}\<%
\\
\>[10][@{}l@{\AgdaIndent{0}}]%
\>[12]\AgdaSymbol{(}\AgdaFunction{substP}\AgdaSpace{}%
\AgdaSymbol{(λ}\AgdaSpace{}%
\AgdaBound{x}\AgdaSpace{}%
\AgdaSymbol{→}\AgdaSpace{}%
\AgdaField{Pf}\AgdaSpace{}%
\AgdaSymbol{(}\AgdaField{◇}\AgdaSpace{}%
\AgdaOperator{\AgdaField{▹ₚ}}\AgdaSpace{}%
\AgdaField{Forall}\AgdaSpace{}%
\AgdaSymbol{(}\AgdaBound{P}\AgdaSpace{}%
\AgdaOperator{\AgdaField{∧}}\AgdaSpace{}%
\AgdaBound{Q}\AgdaSymbol{))}\AgdaSpace{}%
\AgdaSymbol{(}\AgdaField{Forall}\AgdaSpace{}%
\AgdaBound{x}\AgdaSymbol{))}\AgdaSpace{}%
\AgdaField{∧[]}\<%
\\
\>[12][@{}l@{\AgdaIndent{0}}]%
\>[14]\AgdaSymbol{(}\AgdaFunction{substP}\AgdaSpace{}%
\AgdaSymbol{(λ}\AgdaSpace{}%
\AgdaBound{x}\AgdaSpace{}%
\AgdaSymbol{→}\AgdaSpace{}%
\AgdaField{Pf}\AgdaSpace{}%
\AgdaSymbol{(}\AgdaField{◇}\AgdaSpace{}%
\AgdaOperator{\AgdaField{▹ₚ}}\AgdaSpace{}%
\AgdaField{Forall}\AgdaSpace{}%
\AgdaSymbol{(}\AgdaBound{P}\AgdaSpace{}%
\AgdaOperator{\AgdaField{∧}}\AgdaSpace{}%
\AgdaBound{Q}\AgdaSymbol{))}\AgdaSpace{}%
\AgdaBound{x}\AgdaSymbol{)}\AgdaSpace{}%
\AgdaField{Forall[]}\AgdaSpace{}%
\AgdaField{qₚ}\AgdaSymbol{))))))}\<%
\\
%
\>[8]\AgdaSymbol{(}\AgdaFunction{substP}\AgdaSpace{}%
\AgdaSymbol{(λ}\AgdaSpace{}%
\AgdaBound{x}\AgdaSpace{}%
\AgdaSymbol{→}\AgdaSpace{}%
\AgdaField{Pf}\AgdaSpace{}%
\AgdaSymbol{(}\AgdaField{◇}\AgdaSpace{}%
\AgdaOperator{\AgdaField{▹ₚ}}\AgdaSpace{}%
\AgdaField{Forall}\AgdaSpace{}%
\AgdaSymbol{(}\AgdaBound{P}\AgdaSpace{}%
\AgdaOperator{\AgdaField{∧}}\AgdaSpace{}%
\AgdaBound{Q}\AgdaSymbol{))}\AgdaSpace{}%
\AgdaBound{x}\AgdaSymbol{)}\AgdaSpace{}%
\AgdaSymbol{(}\AgdaField{Forall[]}\AgdaSpace{}%
\AgdaOperator{\AgdaFunction{⁻¹}}\AgdaSymbol{)}\<%
\\
\>[8][@{}l@{\AgdaIndent{0}}]%
\>[10]\AgdaSymbol{(}\AgdaField{∀in}\AgdaSpace{}%
\AgdaSymbol{(}\AgdaField{∧out₂}\AgdaSpace{}%
\AgdaSymbol{(}\AgdaField{∀out}\<%
\\
\>[10][@{}l@{\AgdaIndent{0}}]%
\>[12]\AgdaSymbol{(}\AgdaFunction{substP}\AgdaSpace{}%
\AgdaSymbol{(λ}\AgdaSpace{}%
\AgdaBound{x}\AgdaSpace{}%
\AgdaSymbol{→}\AgdaSpace{}%
\AgdaField{Pf}\AgdaSpace{}%
\AgdaSymbol{(}\AgdaField{◇}\AgdaSpace{}%
\AgdaOperator{\AgdaField{▹ₚ}}\AgdaSpace{}%
\AgdaField{Forall}\AgdaSpace{}%
\AgdaSymbol{(}\AgdaBound{P}\AgdaSpace{}%
\AgdaOperator{\AgdaField{∧}}\AgdaSpace{}%
\AgdaBound{Q}\AgdaSymbol{))}\AgdaSpace{}%
\AgdaSymbol{(}\AgdaField{Forall}\AgdaSpace{}%
\AgdaBound{x}\AgdaSymbol{))}\AgdaSpace{}%
\AgdaField{∧[]}\<%
\\
\>[12][@{}l@{\AgdaIndent{0}}]%
\>[14]\AgdaSymbol{(}\AgdaFunction{substP}\AgdaSpace{}%
\AgdaSymbol{(λ}\AgdaSpace{}%
\AgdaBound{x}\AgdaSpace{}%
\AgdaSymbol{→}\AgdaSpace{}%
\AgdaField{Pf}\AgdaSpace{}%
\AgdaSymbol{(}\AgdaField{◇}\AgdaSpace{}%
\AgdaOperator{\AgdaField{▹ₚ}}\AgdaSpace{}%
\AgdaField{Forall}\AgdaSpace{}%
\AgdaSymbol{(}\AgdaBound{P}\AgdaSpace{}%
\AgdaOperator{\AgdaField{∧}}\AgdaSpace{}%
\AgdaBound{Q}\AgdaSymbol{))}\AgdaSpace{}%
\AgdaBound{x}\AgdaSymbol{)}\AgdaSpace{}%
\AgdaField{Forall[]}\AgdaSpace{}%
\AgdaField{qₚ}\AgdaSymbol{))))))))}\<%
\end{code}}

\begin{code}[hide]%
\>[0]\AgdaKeyword{module}\AgdaSpace{}%
\AgdaModule{workInSyntax}\AgdaSpace{}%
\AgdaKeyword{where}\<%
\\
\>[0][@{}l@{\AgdaIndent{0}}]%
\>[2]\AgdaKeyword{import}\AgdaSpace{}%
\AgdaModule{Syntax}\AgdaSpace{}%
\AgdaDatatype{funar}\AgdaSpace{}%
\AgdaDatatype{relar}\ as\ \AgdaModule{S}\<%
\\
%
\>[2]\AgdaKeyword{open}\AgdaSpace{}%
\AgdaModule{Model}\AgdaSpace{}%
\AgdaFunction{S.I}\<%
\\
%
\\[\AgdaEmptyExtraSkip]%
%
\>[2]\AgdaFunction{zero'}\AgdaSpace{}%
\AgdaSymbol{:}\AgdaSpace{}%
\AgdaSymbol{∀\{}\AgdaBound{Γ}\AgdaSymbol{\}}\AgdaSpace{}%
\AgdaSymbol{→}\AgdaSpace{}%
\AgdaFunction{Tm}\AgdaSpace{}%
\AgdaBound{Γ}\<%
\\
%
\>[2]\AgdaFunction{zero'}\AgdaSpace{}%
\AgdaSymbol{\{}\AgdaBound{Γ}\AgdaSymbol{\}}\AgdaSpace{}%
\AgdaSymbol{=}\AgdaSpace{}%
\AgdaFunction{fun}\AgdaSpace{}%
\AgdaSymbol{\{}\AgdaBound{Γ}\AgdaSymbol{\}}\AgdaSpace{}%
\AgdaInductiveConstructor{zero}\AgdaSpace{}%
\AgdaSymbol{\AgdaUnderscore{}}\<%
\\
%
\\[\AgdaEmptyExtraSkip]%
%
\>[2]\AgdaFunction{suc'}\AgdaSpace{}%
\AgdaSymbol{:}\AgdaSpace{}%
\AgdaSymbol{∀\{}\AgdaBound{Γ}\AgdaSymbol{\}}\AgdaSpace{}%
\AgdaSymbol{→}\AgdaSpace{}%
\AgdaFunction{Tm}\AgdaSpace{}%
\AgdaBound{Γ}\AgdaSpace{}%
\AgdaSymbol{→}\AgdaSpace{}%
\AgdaFunction{Tm}\AgdaSpace{}%
\AgdaBound{Γ}\<%
\\
%
\>[2]\AgdaFunction{suc'}\AgdaSpace{}%
\AgdaSymbol{\{}\AgdaBound{Γ}\AgdaSymbol{\}}\AgdaSpace{}%
\AgdaBound{n}\AgdaSpace{}%
\AgdaSymbol{=}\AgdaSpace{}%
\AgdaFunction{fun}\AgdaSpace{}%
\AgdaSymbol{\{}\AgdaBound{Γ}\AgdaSymbol{\}}\AgdaSpace{}%
\AgdaInductiveConstructor{suc}\AgdaSpace{}%
\AgdaSymbol{((}\AgdaBound{n}\AgdaSpace{}%
\AgdaOperator{\AgdaInductiveConstructor{,}}\AgdaSpace{}%
\AgdaSymbol{\AgdaUnderscore{}))}\<%
\\
%
\\[\AgdaEmptyExtraSkip]%
%
\>[2]\AgdaFunction{three}\AgdaSpace{}%
\AgdaSymbol{:}\AgdaSpace{}%
\AgdaSymbol{∀\{}\AgdaBound{Γ}\AgdaSymbol{\}}\AgdaSpace{}%
\AgdaSymbol{→}\AgdaSpace{}%
\AgdaFunction{Tm}\AgdaSpace{}%
\AgdaBound{Γ}\<%
\\
%
\>[2]\AgdaFunction{three}\AgdaSpace{}%
\AgdaSymbol{\{}\AgdaBound{Γ}\AgdaSymbol{\}}\AgdaSpace{}%
\AgdaSymbol{=}\AgdaSpace{}%
\AgdaFunction{suc'}\AgdaSpace{}%
\AgdaSymbol{\{}\AgdaBound{Γ}\AgdaSymbol{\}}\AgdaSpace{}%
\AgdaSymbol{(}\AgdaFunction{suc'}\AgdaSpace{}%
\AgdaSymbol{\{}\AgdaBound{Γ}\AgdaSymbol{\}}\AgdaSpace{}%
\AgdaSymbol{(}\AgdaFunction{suc'}\AgdaSpace{}%
\AgdaSymbol{\{}\AgdaBound{Γ}\AgdaSymbol{\}}\AgdaSpace{}%
\AgdaSymbol{(}\AgdaFunction{zero'}\AgdaSpace{}%
\AgdaSymbol{\{}\AgdaBound{Γ}\AgdaSymbol{\})))}\<%
\\
%
\\[\AgdaEmptyExtraSkip]%
%
\>[2]\AgdaFunction{Eq}\AgdaSpace{}%
\AgdaSymbol{:}\AgdaSpace{}%
\AgdaSymbol{∀\{}\AgdaBound{Γ}\AgdaSymbol{\}}\AgdaSpace{}%
\AgdaSymbol{→}\AgdaSpace{}%
\AgdaFunction{Tm}\AgdaSpace{}%
\AgdaBound{Γ}\AgdaSpace{}%
\AgdaSymbol{→}\AgdaSpace{}%
\AgdaFunction{Tm}\AgdaSpace{}%
\AgdaBound{Γ}\AgdaSpace{}%
\AgdaSymbol{→}\AgdaSpace{}%
\AgdaFunction{For}\AgdaSpace{}%
\AgdaBound{Γ}\<%
\\
%
\>[2]\AgdaFunction{Eq}\AgdaSpace{}%
\AgdaSymbol{\{}\AgdaBound{Γ}\AgdaSymbol{\}}\AgdaSpace{}%
\AgdaBound{u}\AgdaSpace{}%
\AgdaBound{v}\AgdaSpace{}%
\AgdaSymbol{=}\AgdaSpace{}%
\AgdaFunction{Rel}\AgdaSpace{}%
\AgdaSymbol{\{}\AgdaBound{Γ}\AgdaSymbol{\}}\AgdaSpace{}%
\AgdaInductiveConstructor{eq}\AgdaSpace{}%
\AgdaSymbol{(}\AgdaBound{u}\AgdaSpace{}%
\AgdaOperator{\AgdaInductiveConstructor{,}}\AgdaSpace{}%
\AgdaSymbol{(}\AgdaBound{v}\AgdaSpace{}%
\AgdaOperator{\AgdaInductiveConstructor{,}}\AgdaSpace{}%
\AgdaSymbol{\AgdaUnderscore{}))}\<%
\\
%
\\[\AgdaEmptyExtraSkip]%
%
\>[2]\AgdaFunction{Γ₀₁}\AgdaSpace{}%
\AgdaSymbol{=}\AgdaSpace{}%
\AgdaFunction{◇}\AgdaSpace{}%
\AgdaOperator{\AgdaFunction{▹ₚ}}\AgdaSpace{}%
\AgdaSymbol{(}\AgdaFunction{Forall}\AgdaSpace{}%
\AgdaSymbol{\{}\AgdaFunction{◇}\AgdaSymbol{\}}\AgdaSpace{}%
\AgdaSymbol{(}\AgdaFunction{Forall}\AgdaSpace{}%
\AgdaSymbol{\{}\AgdaFunction{◇}\AgdaSpace{}%
\AgdaOperator{\AgdaFunction{▹ₜ}}\AgdaSymbol{\}}\AgdaSpace{}%
\AgdaSymbol{(}\AgdaOperator{\AgdaFunction{\AgdaUnderscore{}⊃\AgdaUnderscore{}}}\AgdaSpace{}%
\AgdaSymbol{\{}\AgdaFunction{◇}\AgdaSpace{}%
\AgdaOperator{\AgdaFunction{▹ₜ}}\AgdaSpace{}%
\AgdaOperator{\AgdaFunction{▹ₜ}}\AgdaSymbol{\}}\AgdaSpace{}%
\AgdaSymbol{(}\AgdaFunction{Eq}\AgdaSpace{}%
\AgdaSymbol{\{}\AgdaFunction{◇}\AgdaSpace{}%
\AgdaOperator{\AgdaFunction{▹ₜ}}\AgdaSpace{}%
\AgdaOperator{\AgdaFunction{▹ₜ}}\AgdaSymbol{\}}\AgdaSpace{}%
\AgdaSymbol{(}\AgdaFunction{qₜ}\AgdaSpace{}%
\AgdaSymbol{\{}\AgdaFunction{◇}\AgdaSymbol{\}}\AgdaSpace{}%
\AgdaOperator{\AgdaFunction{[}}\AgdaSpace{}%
\AgdaFunction{pₜ}\AgdaSpace{}%
\AgdaSymbol{\{}\AgdaFunction{◇}\AgdaSpace{}%
\AgdaOperator{\AgdaFunction{▹ₜ}}\AgdaSymbol{\}}\AgdaSpace{}%
\AgdaOperator{\AgdaFunction{]ᵗ}}\AgdaSymbol{)}\AgdaSpace{}%
\AgdaSymbol{(}\AgdaFunction{qₜ}\AgdaSpace{}%
\AgdaSymbol{\{}\AgdaFunction{◇}\AgdaSpace{}%
\AgdaOperator{\AgdaFunction{▹ₜ}}\AgdaSymbol{\})}\AgdaSpace{}%
\AgdaSymbol{)}\AgdaSpace{}%
\AgdaSymbol{(}\AgdaFunction{Eq}\AgdaSpace{}%
\AgdaSymbol{\{}\AgdaFunction{◇}\AgdaSpace{}%
\AgdaOperator{\AgdaFunction{▹ₜ}}\AgdaSpace{}%
\AgdaOperator{\AgdaFunction{▹ₜ}}\AgdaSymbol{\}}\AgdaSpace{}%
\AgdaSymbol{(}\AgdaFunction{qₜ}\AgdaSpace{}%
\AgdaSymbol{\{}\AgdaFunction{◇}\AgdaSpace{}%
\AgdaOperator{\AgdaFunction{▹ₜ}}\AgdaSymbol{\})}\AgdaSpace{}%
\AgdaSymbol{(}\AgdaFunction{qₜ}\AgdaSpace{}%
\AgdaSymbol{\{}\AgdaFunction{◇}\AgdaSymbol{\}}\AgdaSpace{}%
\AgdaOperator{\AgdaFunction{[}}\AgdaSpace{}%
\AgdaFunction{pₜ}\AgdaSpace{}%
\AgdaSymbol{\{}\AgdaFunction{◇}\AgdaSpace{}%
\AgdaOperator{\AgdaFunction{▹ₜ}}\AgdaSymbol{\}}\AgdaSpace{}%
\AgdaOperator{\AgdaFunction{]ᵗ}}\AgdaSymbol{)))))}\<%
\\
\>[0]\<%
\\
%
\>[2]\AgdaFunction{Γ₀₂}\AgdaSpace{}%
\AgdaSymbol{=}\AgdaSpace{}%
\AgdaFunction{Γ₀₁}\AgdaSpace{}%
\AgdaOperator{\AgdaFunction{▹ₚ}}\AgdaSpace{}%
\AgdaSymbol{(}\AgdaFunction{Forall}\AgdaSpace{}%
\AgdaSymbol{\{}\AgdaFunction{◇}\AgdaSymbol{\}}\AgdaSpace{}%
\AgdaSymbol{(}\AgdaFunction{Forall}\AgdaSpace{}%
\AgdaSymbol{\{}\AgdaFunction{◇}\AgdaSpace{}%
\AgdaOperator{\AgdaFunction{▹ₜ}}\AgdaSymbol{\}}\AgdaSpace{}%
\AgdaSymbol{(}\AgdaFunction{Forall}\AgdaSpace{}%
\AgdaSymbol{\{}\AgdaFunction{◇}\AgdaSpace{}%
\AgdaOperator{\AgdaFunction{▹ₜ}}\AgdaSpace{}%
\AgdaOperator{\AgdaFunction{▹ₜ}}\AgdaSymbol{\}}\AgdaSpace{}%
\AgdaSymbol{(}\AgdaOperator{\AgdaFunction{\AgdaUnderscore{}⊃\AgdaUnderscore{}}}\AgdaSpace{}%
\AgdaSymbol{\{}\AgdaFunction{◇}\AgdaSpace{}%
\AgdaOperator{\AgdaFunction{▹ₜ}}\AgdaSpace{}%
\AgdaOperator{\AgdaFunction{▹ₜ}}\AgdaSpace{}%
\AgdaOperator{\AgdaFunction{▹ₜ}}\AgdaSymbol{\}}\AgdaSpace{}%
\AgdaSymbol{(}\AgdaFunction{Eq}\AgdaSpace{}%
\AgdaSymbol{\{}\AgdaFunction{◇}\AgdaSpace{}%
\AgdaOperator{\AgdaFunction{▹ₜ}}\AgdaSpace{}%
\AgdaOperator{\AgdaFunction{▹ₜ}}\AgdaSpace{}%
\AgdaOperator{\AgdaFunction{▹ₜ}}\AgdaSymbol{\}}\AgdaSpace{}%
\AgdaSymbol{(}\AgdaFunction{qₜ}\AgdaSpace{}%
\AgdaSymbol{\{}\AgdaFunction{◇}\AgdaSymbol{\}}\AgdaSpace{}%
\AgdaOperator{\AgdaFunction{[}}\AgdaSpace{}%
\AgdaFunction{pₜ}\AgdaSpace{}%
\AgdaSymbol{\{}\AgdaFunction{◇}\AgdaSpace{}%
\AgdaOperator{\AgdaFunction{▹ₜ}}\AgdaSymbol{\}}\AgdaSpace{}%
\AgdaOperator{\AgdaFunction{]ᵗ}}\AgdaSpace{}%
\AgdaOperator{\AgdaFunction{[}}\AgdaSpace{}%
\AgdaFunction{pₜ}\AgdaSpace{}%
\AgdaSymbol{\{}\AgdaFunction{◇}\AgdaSpace{}%
\AgdaOperator{\AgdaFunction{▹ₜ}}\AgdaSpace{}%
\AgdaOperator{\AgdaFunction{▹ₜ}}\AgdaSymbol{\}}\AgdaSpace{}%
\AgdaOperator{\AgdaFunction{]ᵗ}}\AgdaSymbol{)}\AgdaSpace{}%
\AgdaSymbol{(}\AgdaFunction{qₜ}\AgdaSpace{}%
\AgdaSymbol{\{}\AgdaFunction{◇}\AgdaSpace{}%
\AgdaOperator{\AgdaFunction{▹ₜ}}\AgdaSymbol{\}}\AgdaSpace{}%
\AgdaOperator{\AgdaFunction{[}}\AgdaSpace{}%
\AgdaFunction{pₜ}\AgdaSpace{}%
\AgdaSymbol{\{}\AgdaFunction{◇}\AgdaSpace{}%
\AgdaOperator{\AgdaFunction{▹ₜ}}\AgdaSpace{}%
\AgdaOperator{\AgdaFunction{▹ₜ}}\AgdaSymbol{\}}\AgdaSpace{}%
\AgdaOperator{\AgdaFunction{]ᵗ}}\AgdaSymbol{))}\AgdaSpace{}%
\AgdaSymbol{(}\AgdaOperator{\AgdaFunction{\AgdaUnderscore{}⊃\AgdaUnderscore{}}}\AgdaSpace{}%
\AgdaSymbol{\{}\AgdaFunction{◇}\AgdaSpace{}%
\AgdaOperator{\AgdaFunction{▹ₜ}}\AgdaSpace{}%
\AgdaOperator{\AgdaFunction{▹ₜ}}\AgdaSpace{}%
\AgdaOperator{\AgdaFunction{▹ₜ}}\AgdaSymbol{\}}\AgdaSpace{}%
\AgdaSymbol{(}\AgdaFunction{Eq}\AgdaSpace{}%
\AgdaSymbol{\{}\AgdaFunction{◇}\AgdaSpace{}%
\AgdaOperator{\AgdaFunction{▹ₜ}}\AgdaSpace{}%
\AgdaOperator{\AgdaFunction{▹ₜ}}\AgdaSpace{}%
\AgdaOperator{\AgdaFunction{▹ₜ}}\AgdaSymbol{\}}\AgdaSpace{}%
\AgdaSymbol{(}\AgdaFunction{qₜ}\AgdaSpace{}%
\AgdaSymbol{\{}\AgdaFunction{◇}\AgdaSpace{}%
\AgdaOperator{\AgdaFunction{▹ₜ}}\AgdaSymbol{\}}\AgdaSpace{}%
\AgdaOperator{\AgdaFunction{[}}\AgdaSpace{}%
\AgdaFunction{pₜ}\AgdaSpace{}%
\AgdaSymbol{\{}\AgdaFunction{◇}\AgdaSpace{}%
\AgdaOperator{\AgdaFunction{▹ₜ}}\AgdaSpace{}%
\AgdaOperator{\AgdaFunction{▹ₜ}}\AgdaSymbol{\}}\AgdaSpace{}%
\AgdaOperator{\AgdaFunction{]ᵗ}}\AgdaSymbol{)}\AgdaSpace{}%
\AgdaSymbol{(}\AgdaFunction{qₜ}\AgdaSpace{}%
\AgdaSymbol{\{}\AgdaFunction{◇}\AgdaSpace{}%
\AgdaOperator{\AgdaFunction{▹ₜ}}\AgdaSpace{}%
\AgdaOperator{\AgdaFunction{▹ₜ}}\AgdaSymbol{\}))}\AgdaSpace{}%
\AgdaSymbol{(}\AgdaFunction{Eq}\AgdaSpace{}%
\AgdaSymbol{\{}\AgdaFunction{◇}\AgdaSpace{}%
\AgdaOperator{\AgdaFunction{▹ₜ}}\AgdaSpace{}%
\AgdaOperator{\AgdaFunction{▹ₜ}}\AgdaSpace{}%
\AgdaOperator{\AgdaFunction{▹ₜ}}\AgdaSymbol{\}}\AgdaSpace{}%
\AgdaSymbol{(}\AgdaFunction{qₜ}\AgdaSpace{}%
\AgdaSymbol{\{}\AgdaFunction{◇}\AgdaSymbol{\}}\AgdaSpace{}%
\AgdaOperator{\AgdaFunction{[}}\AgdaSpace{}%
\AgdaFunction{pₜ}\AgdaSpace{}%
\AgdaSymbol{\{}\AgdaFunction{◇}\AgdaSpace{}%
\AgdaOperator{\AgdaFunction{▹ₜ}}\AgdaSymbol{\}}\AgdaSpace{}%
\AgdaOperator{\AgdaFunction{]ᵗ}}\AgdaSpace{}%
\AgdaOperator{\AgdaFunction{[}}\AgdaSpace{}%
\AgdaFunction{pₜ}\AgdaSpace{}%
\AgdaSymbol{\{}\AgdaFunction{◇}\AgdaSpace{}%
\AgdaOperator{\AgdaFunction{▹ₜ}}\AgdaSpace{}%
\AgdaOperator{\AgdaFunction{▹ₜ}}\AgdaSymbol{\}}\AgdaSpace{}%
\AgdaOperator{\AgdaFunction{]ᵗ}}\AgdaSymbol{)}\AgdaSpace{}%
\AgdaSymbol{(}\AgdaFunction{qₜ}\AgdaSpace{}%
\AgdaSymbol{\{}\AgdaFunction{◇}\AgdaSpace{}%
\AgdaOperator{\AgdaFunction{▹ₜ}}\AgdaSpace{}%
\AgdaOperator{\AgdaFunction{▹ₜ}}\AgdaSymbol{\})))))))}\<%
\\
\>[0]\<%
\\
%
\>[2]\AgdaFunction{Γ}\AgdaSpace{}%
\AgdaSymbol{=}\AgdaSpace{}%
\AgdaFunction{Γ₀₂}\AgdaSpace{}%
\AgdaOperator{\AgdaFunction{▹ₜ}}\AgdaSpace{}%
\AgdaOperator{\AgdaFunction{▹ₚ}}\AgdaSpace{}%
\AgdaFunction{Eq}\AgdaSpace{}%
\AgdaSymbol{\{}\AgdaFunction{◇}\AgdaSpace{}%
\AgdaOperator{\AgdaFunction{▹ₜ}}\AgdaSymbol{\}}\AgdaSpace{}%
\AgdaSymbol{(}\AgdaFunction{qₜ}\AgdaSpace{}%
\AgdaSymbol{\{}\AgdaFunction{◇}\AgdaSymbol{\})}\AgdaSpace{}%
\AgdaSymbol{(}\AgdaFunction{three}\AgdaSpace{}%
\AgdaSymbol{\{}\AgdaFunction{◇}\AgdaSpace{}%
\AgdaOperator{\AgdaFunction{▹ₜ}}\AgdaSymbol{\})}\<%
\\
%
\\[\AgdaEmptyExtraSkip]%
%
\>[2]\AgdaFunction{A⊃A}\AgdaSpace{}%
\AgdaSymbol{:}\AgdaSpace{}%
\AgdaSymbol{\{}\AgdaBound{A}\AgdaSpace{}%
\AgdaSymbol{:}\AgdaSpace{}%
\AgdaFunction{For}\AgdaSpace{}%
\AgdaFunction{◇}\AgdaSymbol{\}}\AgdaSpace{}%
\AgdaSymbol{→}\AgdaSpace{}%
\AgdaFunction{Pf}\AgdaSpace{}%
\AgdaFunction{◇}\AgdaSpace{}%
\AgdaSymbol{(}\AgdaOperator{\AgdaFunction{\AgdaUnderscore{}⊃\AgdaUnderscore{}}}\AgdaSpace{}%
\AgdaSymbol{\{}\AgdaFunction{◇}\AgdaSymbol{\}}\AgdaSpace{}%
\AgdaBound{A}\AgdaSpace{}%
\AgdaBound{A}\AgdaSymbol{)}\<%
\\
%
\>[2]\AgdaFunction{A⊃A}\AgdaSpace{}%
\AgdaSymbol{=}\AgdaSpace{}%
\AgdaFunction{⊃in}\AgdaSpace{}%
\AgdaFunction{qₚ}\<%
\end{code}
\newcommand{\pSynE}{
\begin{code}%
%
\>[2]\AgdaFunction{A⊃B⊃A}\AgdaSpace{}%
\AgdaSymbol{:}\AgdaSpace{}%
\AgdaSymbol{\{}\AgdaBound{A}\AgdaSpace{}%
\AgdaBound{B}\AgdaSpace{}%
\AgdaSymbol{:}\AgdaSpace{}%
\AgdaFunction{For}\AgdaSpace{}%
\AgdaFunction{◇}\AgdaSymbol{\}}\AgdaSpace{}%
\AgdaSymbol{→}\AgdaSpace{}%
\AgdaFunction{Pf}\AgdaSpace{}%
\AgdaFunction{◇}\AgdaSpace{}%
\AgdaSymbol{(}\AgdaOperator{\AgdaFunction{\AgdaUnderscore{}⊃\AgdaUnderscore{}}}\AgdaSpace{}%
\AgdaSymbol{\{}\AgdaFunction{◇}\AgdaSymbol{\}}\AgdaSpace{}%
\AgdaBound{A}\AgdaSpace{}%
\AgdaSymbol{(}\AgdaOperator{\AgdaFunction{\AgdaUnderscore{}⊃\AgdaUnderscore{}}}\AgdaSpace{}%
\AgdaSymbol{\{}\AgdaFunction{◇}\AgdaSymbol{\}}\AgdaSpace{}%
\AgdaBound{B}\AgdaSpace{}%
\AgdaBound{A}\AgdaSymbol{))}\<%
\\
%
\>[2]\AgdaFunction{A⊃B⊃A}\AgdaSpace{}%
\AgdaSymbol{=}\AgdaSpace{}%
\AgdaFunction{⊃in}\AgdaSpace{}%
\AgdaSymbol{(}\AgdaFunction{⊃in}\AgdaSpace{}%
\AgdaSymbol{(}\AgdaFunction{qₚ}\AgdaSpace{}%
\AgdaOperator{\AgdaFunction{[}}\AgdaSpace{}%
\AgdaFunction{pₚ}\AgdaSpace{}%
\AgdaOperator{\AgdaFunction{]ᵖ}}\AgdaSymbol{))}\<%
\end{code}}
\begin{code}[hide]%
%
\>[2]\AgdaFunction{A∧B⊃B∨A}\AgdaSpace{}%
\AgdaSymbol{:}\AgdaSpace{}%
\AgdaSymbol{\{}\AgdaBound{A}\AgdaSpace{}%
\AgdaBound{B}\AgdaSpace{}%
\AgdaSymbol{:}\AgdaSpace{}%
\AgdaFunction{For}\AgdaSpace{}%
\AgdaFunction{◇}\AgdaSymbol{\}}\AgdaSpace{}%
\AgdaSymbol{→}\AgdaSpace{}%
\AgdaFunction{Pf}\AgdaSpace{}%
\AgdaSymbol{(}\AgdaFunction{◇}\AgdaSymbol{)}\AgdaSpace{}%
\AgdaSymbol{(}\AgdaOperator{\AgdaFunction{\AgdaUnderscore{}⊃\AgdaUnderscore{}}}\AgdaSpace{}%
\AgdaSymbol{\{}\AgdaFunction{◇}\AgdaSymbol{\}}\AgdaSpace{}%
\AgdaSymbol{(}\AgdaOperator{\AgdaFunction{\AgdaUnderscore{}∧\AgdaUnderscore{}}}\AgdaSpace{}%
\AgdaSymbol{\{}\AgdaFunction{◇}\AgdaSymbol{\}}\AgdaSpace{}%
\AgdaBound{A}\AgdaSpace{}%
\AgdaBound{B}\AgdaSymbol{)}\AgdaSpace{}%
\AgdaSymbol{(}\AgdaOperator{\AgdaFunction{\AgdaUnderscore{}∨\AgdaUnderscore{}}}\AgdaSpace{}%
\AgdaSymbol{\{}\AgdaFunction{◇}\AgdaSymbol{\}}\AgdaSpace{}%
\AgdaBound{A}\AgdaSpace{}%
\AgdaBound{B}\AgdaSymbol{)}\AgdaSpace{}%
\AgdaSymbol{)}\<%
\\
%
\>[2]\AgdaFunction{A∧B⊃B∨A}\AgdaSpace{}%
\AgdaSymbol{=}\AgdaSpace{}%
\AgdaFunction{⊃in}\AgdaSpace{}%
\AgdaSymbol{(}\AgdaFunction{∨in₁}\AgdaSpace{}%
\AgdaSymbol{(}\AgdaFunction{∧out₁}\AgdaSpace{}%
\AgdaFunction{qₚ}\AgdaSymbol{))}\<%
\end{code}
\newcommand{\pSynC}{
\begin{code}%
%
\>[2]\AgdaFunction{∀P∧Q⊃∀P∧∀Q}\AgdaSpace{}%
\AgdaSymbol{:}\AgdaSpace{}%
\AgdaSymbol{\{}\AgdaBound{P}\AgdaSpace{}%
\AgdaBound{Q}\AgdaSpace{}%
\AgdaSymbol{:}\AgdaSpace{}%
\AgdaFunction{For}\AgdaSpace{}%
\AgdaSymbol{(}\AgdaFunction{◇}\AgdaSpace{}%
\AgdaOperator{\AgdaFunction{▹ₜ}}\AgdaSymbol{)\}}\AgdaSpace{}%
\AgdaSymbol{→}\AgdaSpace{}%
\AgdaFunction{Pf}\AgdaSpace{}%
\AgdaSymbol{(}\AgdaFunction{◇}\AgdaSymbol{)}\AgdaSpace{}%
\AgdaSymbol{(}\AgdaSpace{}%
\AgdaOperator{\AgdaFunction{\AgdaUnderscore{}⊃\AgdaUnderscore{}}}\AgdaSpace{}%
\AgdaSymbol{\{}\AgdaFunction{◇}\AgdaSymbol{\}}\<%
\\
\>[2][@{}l@{\AgdaIndent{0}}]%
\>[8]\AgdaSymbol{(}\AgdaFunction{Forall}\AgdaSpace{}%
\AgdaSymbol{\{}\AgdaFunction{◇}\AgdaSymbol{\}}\AgdaSpace{}%
\AgdaSymbol{(}\AgdaOperator{\AgdaFunction{\AgdaUnderscore{}∧\AgdaUnderscore{}}}\AgdaSpace{}%
\AgdaSymbol{\{}\AgdaFunction{◇}\AgdaSpace{}%
\AgdaOperator{\AgdaFunction{▹ₜ}}\AgdaSymbol{\}}\AgdaSpace{}%
\AgdaBound{P}\AgdaSpace{}%
\AgdaBound{Q}\AgdaSymbol{))}\<%
\\
%
\>[8]\AgdaSymbol{(}\AgdaOperator{\AgdaFunction{\AgdaUnderscore{}∧\AgdaUnderscore{}}}\AgdaSpace{}%
\AgdaSymbol{\{}\AgdaFunction{◇}\AgdaSymbol{\}}\AgdaSpace{}%
\AgdaSymbol{(}\AgdaFunction{Forall}\AgdaSpace{}%
\AgdaSymbol{\{}\AgdaFunction{◇}\AgdaSymbol{\}}\AgdaSpace{}%
\AgdaBound{P}\AgdaSymbol{)}\AgdaSpace{}%
\AgdaSymbol{(}\AgdaFunction{Forall}\AgdaSpace{}%
\AgdaSymbol{\{}\AgdaFunction{◇}\AgdaSymbol{\}}\AgdaSpace{}%
\AgdaBound{Q}\AgdaSymbol{)))}\<%
\\
%
\>[2]\AgdaFunction{∀P∧Q⊃∀P∧∀Q}\AgdaSpace{}%
\AgdaSymbol{=}\AgdaSpace{}%
\AgdaFunction{⊃in}\AgdaSpace{}%
\AgdaSymbol{(}\AgdaFunction{∧in}\AgdaSpace{}%
\AgdaSymbol{(}\AgdaFunction{∀in}\AgdaSpace{}%
\AgdaSymbol{(}\AgdaFunction{∧out₁}\AgdaSpace{}%
\AgdaSymbol{(}\AgdaFunction{∀out}\AgdaSpace{}%
\AgdaFunction{qₚ}\AgdaSymbol{)))}\AgdaSpace{}%
\AgdaSymbol{(}\AgdaFunction{∀in}\AgdaSpace{}%
\AgdaSymbol{(}\AgdaFunction{∧out₂}\AgdaSpace{}%
\AgdaSymbol{(}\AgdaFunction{∀out}\AgdaSpace{}%
\AgdaFunction{qₚ}\AgdaSymbol{))))}\<%
\end{code}}

\begin{code}[hide]%
\>[0]\AgdaKeyword{module}\AgdaSpace{}%
\AgdaModule{workInTarski}\AgdaSpace{}%
\AgdaKeyword{where}\<%
\\
%
\\[\AgdaEmptyExtraSkip]%
\>[0][@{}l@{\AgdaIndent{0}}]%
\>[2]\AgdaFunction{ℕfun}\AgdaSpace{}%
\AgdaSymbol{:}\AgdaSpace{}%
\AgdaSymbol{(}\AgdaBound{n}\AgdaSpace{}%
\AgdaSymbol{:}\AgdaSpace{}%
\AgdaDatatype{ℕ}\AgdaSymbol{)}\AgdaSpace{}%
\AgdaSymbol{→}\AgdaSpace{}%
\AgdaDatatype{funar}\AgdaSpace{}%
\AgdaBound{n}\AgdaSpace{}%
\AgdaSymbol{→}\AgdaSpace{}%
\AgdaDatatype{ℕ}\AgdaSpace{}%
\AgdaOperator{\AgdaFunction{\textasciicircum{}}}\AgdaSpace{}%
\AgdaBound{n}\AgdaSpace{}%
\AgdaSymbol{→}\AgdaSpace{}%
\AgdaDatatype{ℕ}\<%
\\
%
\>[2]\AgdaFunction{ℕfun}\AgdaSpace{}%
\AgdaSymbol{\AgdaUnderscore{}}\AgdaSpace{}%
\AgdaInductiveConstructor{zero}\AgdaSpace{}%
\AgdaSymbol{\AgdaUnderscore{}}\AgdaSpace{}%
\AgdaSymbol{=}\AgdaSpace{}%
\AgdaInductiveConstructor{zero}\<%
\\
%
\>[2]\AgdaFunction{ℕfun}\AgdaSpace{}%
\AgdaSymbol{\AgdaUnderscore{}}\AgdaSpace{}%
\AgdaInductiveConstructor{suc}\AgdaSpace{}%
\AgdaSymbol{(}\AgdaBound{n}\AgdaSpace{}%
\AgdaOperator{\AgdaInductiveConstructor{,}}\AgdaSpace{}%
\AgdaSymbol{\AgdaUnderscore{})}\AgdaSpace{}%
\AgdaSymbol{=}\AgdaSpace{}%
\AgdaInductiveConstructor{suc}\AgdaSpace{}%
\AgdaBound{n}\<%
\\
%
\>[2]\AgdaFunction{ℕfun}\AgdaSpace{}%
\AgdaSymbol{\AgdaUnderscore{}}\AgdaSpace{}%
\AgdaInductiveConstructor{+'}\AgdaSpace{}%
\AgdaSymbol{(}\AgdaBound{m}\AgdaSpace{}%
\AgdaOperator{\AgdaInductiveConstructor{,}}\AgdaSpace{}%
\AgdaBound{n}\AgdaSpace{}%
\AgdaOperator{\AgdaInductiveConstructor{,}}\AgdaSpace{}%
\AgdaSymbol{\AgdaUnderscore{})}\AgdaSpace{}%
\AgdaSymbol{=}\AgdaSpace{}%
\AgdaBound{m}\AgdaSpace{}%
\AgdaOperator{\AgdaPrimitive{+}}\AgdaSpace{}%
\AgdaBound{n}\<%
\\
%
\\[\AgdaEmptyExtraSkip]%
%
\>[2]\AgdaFunction{ℕRel}\AgdaSpace{}%
\AgdaSymbol{:}\AgdaSpace{}%
\AgdaSymbol{(}\AgdaBound{n}\AgdaSpace{}%
\AgdaSymbol{:}\AgdaSpace{}%
\AgdaDatatype{ℕ}\AgdaSymbol{)}\AgdaSpace{}%
\AgdaSymbol{→}\AgdaSpace{}%
\AgdaDatatype{relar}\AgdaSpace{}%
\AgdaBound{n}\AgdaSpace{}%
\AgdaSymbol{→}\AgdaSpace{}%
\AgdaDatatype{ℕ}\AgdaSpace{}%
\AgdaOperator{\AgdaFunction{\textasciicircum{}}}\AgdaSpace{}%
\AgdaBound{n}\AgdaSpace{}%
\AgdaSymbol{→}\AgdaSpace{}%
\AgdaPrimitive{Prop}\<%
\\
%
\>[2]\AgdaFunction{ℕRel}\AgdaSpace{}%
\AgdaSymbol{\AgdaUnderscore{}}\AgdaSpace{}%
\AgdaInductiveConstructor{eq}\AgdaSpace{}%
\AgdaSymbol{(}\AgdaBound{m}\AgdaSpace{}%
\AgdaOperator{\AgdaInductiveConstructor{,}}\AgdaSpace{}%
\AgdaBound{n}\AgdaSpace{}%
\AgdaOperator{\AgdaInductiveConstructor{,}}\AgdaSpace{}%
\AgdaSymbol{\AgdaUnderscore{})}\AgdaSpace{}%
\AgdaSymbol{=}\AgdaSpace{}%
\AgdaBound{m}\AgdaSpace{}%
\AgdaOperator{\AgdaDatatype{≡}}\AgdaSpace{}%
\AgdaBound{n}\<%
\\
%
\\[\AgdaEmptyExtraSkip]%
%
\>[2]\AgdaKeyword{import}\AgdaSpace{}%
\AgdaModule{Tarski}\AgdaSpace{}%
\AgdaDatatype{funar}\AgdaSpace{}%
\AgdaDatatype{relar}\AgdaSpace{}%
\AgdaDatatype{ℕ}\AgdaSpace{}%
\AgdaFunction{ℕfun}\AgdaSpace{}%
\AgdaFunction{ℕRel}\ as\ \AgdaModule{T}\<%
\\
%
\>[2]\AgdaKeyword{open}\AgdaSpace{}%
\AgdaModule{Model}\AgdaSpace{}%
\AgdaFunction{T.T}\<%
\\
\>[0]\<%
\\
%
\>[2]\AgdaFunction{zero'}\AgdaSpace{}%
\AgdaSymbol{:}\AgdaSpace{}%
\AgdaSymbol{∀\{}\AgdaBound{Γ}\AgdaSymbol{\}}\AgdaSpace{}%
\AgdaSymbol{→}\AgdaSpace{}%
\AgdaFunction{Tm}\AgdaSpace{}%
\AgdaBound{Γ}\<%
\\
%
\>[2]\AgdaFunction{zero'}\AgdaSpace{}%
\AgdaSymbol{\{}\AgdaBound{Γ}\AgdaSymbol{\}}\AgdaSpace{}%
\AgdaSymbol{=}\AgdaSpace{}%
\AgdaFunction{fun}\AgdaSpace{}%
\AgdaInductiveConstructor{zero}\AgdaSpace{}%
\AgdaSymbol{\AgdaUnderscore{}}\<%
\\
%
\\[\AgdaEmptyExtraSkip]%
%
\>[2]\AgdaFunction{suc'}\AgdaSpace{}%
\AgdaSymbol{:}\AgdaSpace{}%
\AgdaSymbol{∀\{}\AgdaBound{Γ}\AgdaSymbol{\}}\AgdaSpace{}%
\AgdaSymbol{→}\AgdaSpace{}%
\AgdaFunction{Tm}\AgdaSpace{}%
\AgdaBound{Γ}\AgdaSpace{}%
\AgdaSymbol{→}\AgdaSpace{}%
\AgdaFunction{Tm}\AgdaSpace{}%
\AgdaBound{Γ}\<%
\\
%
\>[2]\AgdaFunction{suc'}\AgdaSpace{}%
\AgdaSymbol{\{}\AgdaBound{Γ}\AgdaSymbol{\}}\AgdaSpace{}%
\AgdaBound{n}\AgdaSpace{}%
\AgdaSymbol{=}\AgdaSpace{}%
\AgdaFunction{fun}\AgdaSpace{}%
\AgdaSymbol{\{}\AgdaBound{Γ}\AgdaSymbol{\}}\AgdaSpace{}%
\AgdaInductiveConstructor{suc}\AgdaSpace{}%
\AgdaSymbol{(}\AgdaBound{n}\AgdaSpace{}%
\AgdaOperator{\AgdaInductiveConstructor{,}}\AgdaSpace{}%
\AgdaSymbol{\AgdaUnderscore{})}\<%
\\
%
\\[\AgdaEmptyExtraSkip]%
%
\>[2]\AgdaFunction{three}\AgdaSpace{}%
\AgdaSymbol{:}\AgdaSpace{}%
\AgdaSymbol{∀\{}\AgdaBound{Γ}\AgdaSymbol{\}}\AgdaSpace{}%
\AgdaSymbol{→}\AgdaSpace{}%
\AgdaFunction{Tm}\AgdaSpace{}%
\AgdaBound{Γ}\<%
\\
%
\>[2]\AgdaFunction{three}\AgdaSpace{}%
\AgdaSymbol{\{}\AgdaBound{Γ}\AgdaSymbol{\}}\AgdaSpace{}%
\AgdaSymbol{=}\AgdaSpace{}%
\AgdaFunction{suc'}\AgdaSpace{}%
\AgdaSymbol{\{}\AgdaBound{Γ}\AgdaSymbol{\}}\AgdaSpace{}%
\AgdaSymbol{(}\AgdaFunction{suc'}\AgdaSpace{}%
\AgdaSymbol{\{}\AgdaBound{Γ}\AgdaSymbol{\}}\AgdaSpace{}%
\AgdaSymbol{(}\AgdaFunction{suc'}\AgdaSpace{}%
\AgdaSymbol{\{}\AgdaBound{Γ}\AgdaSymbol{\}}\AgdaSpace{}%
\AgdaSymbol{(}\AgdaFunction{zero'}\AgdaSpace{}%
\AgdaSymbol{\{}\AgdaBound{Γ}\AgdaSymbol{\})))}\<%
\\
%
\\[\AgdaEmptyExtraSkip]%
%
\>[2]\AgdaFunction{Eq}\AgdaSpace{}%
\AgdaSymbol{:}\AgdaSpace{}%
\AgdaSymbol{∀\{}\AgdaBound{Γ}\AgdaSymbol{\}}\AgdaSpace{}%
\AgdaSymbol{→}\AgdaSpace{}%
\AgdaFunction{Tm}\AgdaSpace{}%
\AgdaBound{Γ}\AgdaSpace{}%
\AgdaSymbol{→}\AgdaSpace{}%
\AgdaFunction{Tm}\AgdaSpace{}%
\AgdaBound{Γ}\AgdaSpace{}%
\AgdaSymbol{→}\AgdaSpace{}%
\AgdaFunction{For}\AgdaSpace{}%
\AgdaBound{Γ}\<%
\\
%
\>[2]\AgdaFunction{Eq}\AgdaSpace{}%
\AgdaSymbol{\{}\AgdaBound{Γ}\AgdaSymbol{\}}\AgdaSpace{}%
\AgdaBound{u}\AgdaSpace{}%
\AgdaBound{v}\AgdaSpace{}%
\AgdaSymbol{=}\AgdaSpace{}%
\AgdaFunction{Rel}\AgdaSpace{}%
\AgdaSymbol{\{}\AgdaBound{Γ}\AgdaSymbol{\}}\AgdaSpace{}%
\AgdaInductiveConstructor{eq}\AgdaSpace{}%
\AgdaSymbol{(}\AgdaBound{u}\AgdaSpace{}%
\AgdaOperator{\AgdaInductiveConstructor{,}}\AgdaSpace{}%
\AgdaSymbol{(}\AgdaBound{v}\AgdaSpace{}%
\AgdaOperator{\AgdaInductiveConstructor{,}}\AgdaSpace{}%
\AgdaSymbol{\AgdaUnderscore{}))}\<%
\\
%
\\[\AgdaEmptyExtraSkip]%
%
\>[2]\AgdaFunction{p}\AgdaSpace{}%
\AgdaSymbol{:}\AgdaSpace{}%
\AgdaFunction{Pf}\AgdaSpace{}%
\AgdaSymbol{(}\AgdaFunction{◇}\AgdaSpace{}%
\AgdaOperator{\AgdaFunction{▹ₜ}}\AgdaSpace{}%
\AgdaOperator{\AgdaFunction{▹ₚ}}\AgdaSpace{}%
\AgdaFunction{Eq}\AgdaSpace{}%
\AgdaSymbol{\{}\AgdaFunction{◇}\AgdaSpace{}%
\AgdaOperator{\AgdaFunction{▹ₜ}}\AgdaSymbol{\}}\AgdaSpace{}%
\AgdaSymbol{(}\AgdaFunction{qₜ}\AgdaSpace{}%
\AgdaSymbol{\{}\AgdaFunction{◇}\AgdaSymbol{\})}\AgdaSpace{}%
\AgdaSymbol{(}\AgdaFunction{three}\AgdaSpace{}%
\AgdaSymbol{\{}\AgdaFunction{◇}\AgdaSpace{}%
\AgdaOperator{\AgdaFunction{▹ₜ}}\AgdaSymbol{\}))}\AgdaSpace{}%
\AgdaSymbol{(}\AgdaFunction{Forall}\AgdaSpace{}%
\AgdaSymbol{(}\AgdaFunction{Eq}\AgdaSpace{}%
\AgdaFunction{qₜ}\AgdaSpace{}%
\AgdaFunction{three}\AgdaSpace{}%
\AgdaOperator{\AgdaFunction{⊃}}\AgdaSpace{}%
\AgdaFunction{Eq}\AgdaSpace{}%
\AgdaFunction{qₜ}\AgdaSpace{}%
\AgdaSymbol{(}\AgdaFunction{qₜ}\AgdaSpace{}%
\AgdaOperator{\AgdaFunction{[}}\AgdaSpace{}%
\AgdaFunction{pₚ}\AgdaSpace{}%
\AgdaOperator{\AgdaFunction{]ᵗ}}\AgdaSpace{}%
\AgdaOperator{\AgdaFunction{[}}\AgdaSpace{}%
\AgdaFunction{pₜ}\AgdaSpace{}%
\AgdaOperator{\AgdaFunction{]ᵗ}}\AgdaSymbol{)))}\<%
\\
%
\>[2]\AgdaFunction{p}\AgdaSpace{}%
\AgdaSymbol{=}\AgdaSpace{}%
\AgdaSymbol{λ}\AgdaSpace{}%
\AgdaBound{γ}\AgdaSpace{}%
\AgdaBound{d}\AgdaSpace{}%
\AgdaBound{x}\AgdaSpace{}%
\AgdaSymbol{→}\AgdaSpace{}%
\AgdaSymbol{(}\AgdaField{un}\AgdaSpace{}%
\AgdaSymbol{(}\AgdaField{snd}\AgdaSpace{}%
\AgdaBound{γ}\AgdaSymbol{)}\AgdaSpace{}%
\AgdaOperator{\AgdaFunction{◾}}\AgdaSpace{}%
\AgdaBound{x}\AgdaSpace{}%
\AgdaOperator{\AgdaFunction{⁻¹}}\AgdaSymbol{)}\AgdaSpace{}%
\AgdaOperator{\AgdaFunction{⁻¹}}\<%
\\
%
\\[\AgdaEmptyExtraSkip]%
%
\>[2]\AgdaFunction{A⊃A}\AgdaSpace{}%
\AgdaSymbol{:}\AgdaSpace{}%
\AgdaSymbol{\{}\AgdaBound{A}\AgdaSpace{}%
\AgdaSymbol{:}\AgdaSpace{}%
\AgdaFunction{For}\AgdaSpace{}%
\AgdaFunction{◇}\AgdaSymbol{\}}\AgdaSpace{}%
\AgdaSymbol{→}\AgdaSpace{}%
\AgdaFunction{Pf}\AgdaSpace{}%
\AgdaSymbol{(}\AgdaFunction{◇}\AgdaSymbol{)}\AgdaSpace{}%
\AgdaSymbol{(}\AgdaBound{A}\AgdaSpace{}%
\AgdaOperator{\AgdaFunction{⊃}}\AgdaSpace{}%
\AgdaBound{A}\AgdaSymbol{)}\<%
\\
%
\>[2]\AgdaFunction{A⊃A}\AgdaSpace{}%
\AgdaSymbol{=}\AgdaSpace{}%
\AgdaFunction{⊃in}\AgdaSpace{}%
\AgdaSymbol{λ}\AgdaSpace{}%
\AgdaBound{γ}\AgdaSpace{}%
\AgdaSymbol{→}\AgdaSpace{}%
\AgdaField{un}\AgdaSpace{}%
\AgdaSymbol{(}\AgdaField{snd}\AgdaSpace{}%
\AgdaBound{γ}\AgdaSymbol{)}\<%
\end{code}
\newcommand{\pTarE}{
\begin{code}%
%
\>[2]\AgdaFunction{A⊃B⊃A}\AgdaSpace{}%
\AgdaSymbol{:}\AgdaSpace{}%
\AgdaSymbol{\{}\AgdaBound{A}\AgdaSpace{}%
\AgdaBound{B}\AgdaSpace{}%
\AgdaSymbol{:}\AgdaSpace{}%
\AgdaFunction{For}\AgdaSpace{}%
\AgdaFunction{◇}\AgdaSymbol{\}}\AgdaSpace{}%
\AgdaSymbol{→}\AgdaSpace{}%
\AgdaFunction{Pf}\AgdaSpace{}%
\AgdaSymbol{(}\AgdaFunction{◇}\AgdaSymbol{)}\AgdaSpace{}%
\AgdaSymbol{(}\AgdaBound{A}\AgdaSpace{}%
\AgdaOperator{\AgdaFunction{⊃}}\AgdaSpace{}%
\AgdaBound{B}\AgdaSpace{}%
\AgdaOperator{\AgdaFunction{⊃}}\AgdaSpace{}%
\AgdaBound{A}\AgdaSymbol{)}\<%
\\
%
\>[2]\AgdaFunction{A⊃B⊃A}\AgdaSpace{}%
\AgdaBound{γ}\AgdaSpace{}%
\AgdaBound{a}\AgdaSpace{}%
\AgdaBound{b}\AgdaSpace{}%
\AgdaSymbol{=}\AgdaSpace{}%
\AgdaBound{a}\<%
\end{code}}
\begin{code}[hide]%
%
\>[2]\AgdaFunction{A⊃B⊃A'}\AgdaSpace{}%
\AgdaSymbol{:}\AgdaSpace{}%
\AgdaSymbol{\{}\AgdaBound{A}\AgdaSpace{}%
\AgdaBound{B}\AgdaSpace{}%
\AgdaSymbol{:}\AgdaSpace{}%
\AgdaFunction{For}\AgdaSpace{}%
\AgdaFunction{◇}\AgdaSymbol{\}}\AgdaSpace{}%
\AgdaSymbol{→}\AgdaSpace{}%
\AgdaFunction{Pf}\AgdaSpace{}%
\AgdaSymbol{(}\AgdaFunction{◇}\AgdaSymbol{)}\AgdaSpace{}%
\AgdaSymbol{(}\AgdaBound{A}\AgdaSpace{}%
\AgdaOperator{\AgdaFunction{⊃}}\AgdaSpace{}%
\AgdaBound{B}\AgdaSpace{}%
\AgdaOperator{\AgdaFunction{⊃}}\AgdaSpace{}%
\AgdaBound{A}\AgdaSymbol{)}\<%
\\
%
\>[2]\AgdaFunction{A⊃B⊃A'}\AgdaSpace{}%
\AgdaSymbol{=}\AgdaSpace{}%
\AgdaFunction{⊃in}\AgdaSpace{}%
\AgdaSymbol{λ}\AgdaSpace{}%
\AgdaBound{γ}\AgdaSpace{}%
\AgdaBound{x}\AgdaSpace{}%
\AgdaSymbol{→}\AgdaSpace{}%
\AgdaField{un}\AgdaSpace{}%
\AgdaSymbol{(}\AgdaField{snd}\AgdaSpace{}%
\AgdaBound{γ}\AgdaSymbol{)}\<%
\\
\>[0]\<%
\\
%
\>[2]\AgdaFunction{A∧B⊃B∨A}\AgdaSpace{}%
\AgdaSymbol{:}\AgdaSpace{}%
\AgdaSymbol{\{}\AgdaBound{A}\AgdaSpace{}%
\AgdaBound{B}\AgdaSpace{}%
\AgdaSymbol{:}\AgdaSpace{}%
\AgdaFunction{For}\AgdaSpace{}%
\AgdaFunction{◇}\AgdaSymbol{\}}\AgdaSpace{}%
\AgdaSymbol{→}\AgdaSpace{}%
\AgdaFunction{Pf}\AgdaSpace{}%
\AgdaSymbol{(}\AgdaFunction{◇}\AgdaSymbol{)}\AgdaSpace{}%
\AgdaSymbol{(}\AgdaBound{A}\AgdaSpace{}%
\AgdaOperator{\AgdaFunction{∧}}\AgdaSpace{}%
\AgdaBound{B}\AgdaSpace{}%
\AgdaOperator{\AgdaFunction{⊃}}\AgdaSpace{}%
\AgdaBound{B}\AgdaSpace{}%
\AgdaOperator{\AgdaFunction{∨}}\AgdaSpace{}%
\AgdaBound{A}\AgdaSymbol{)}\<%
\\
%
\>[2]\AgdaFunction{A∧B⊃B∨A}\AgdaSpace{}%
\AgdaSymbol{=}\AgdaSpace{}%
\AgdaFunction{⊃in}\AgdaSpace{}%
\AgdaSymbol{λ}\AgdaSpace{}%
\AgdaBound{γ}\AgdaSpace{}%
\AgdaSymbol{→}\AgdaSpace{}%
\AgdaInductiveConstructor{inl}\AgdaSpace{}%
\AgdaSymbol{(}\AgdaField{snd}\AgdaSpace{}%
\AgdaSymbol{(}\AgdaField{un}\AgdaSpace{}%
\AgdaSymbol{(}\AgdaField{snd}\AgdaSpace{}%
\AgdaBound{γ}\AgdaSymbol{)))}\<%
\end{code}
\newcommand{\pTarC}{
\begin{code}%
%
\>[2]\AgdaFunction{∀P∧Q⊃∀P∧∀Q}\AgdaSpace{}%
\AgdaSymbol{:}%
\>[901I]\AgdaSymbol{\{}\AgdaBound{P}\AgdaSpace{}%
\AgdaBound{Q}\AgdaSpace{}%
\AgdaSymbol{:}\AgdaSpace{}%
\AgdaFunction{For}\AgdaSpace{}%
\AgdaSymbol{(}\AgdaFunction{◇}\AgdaSpace{}%
\AgdaOperator{\AgdaFunction{▹ₜ}}\AgdaSymbol{)\}}\AgdaSpace{}%
\AgdaSymbol{→}\<%
\\
\>[901I][@{}l@{\AgdaIndent{0}}]%
\>[16]\AgdaFunction{Pf}\AgdaSpace{}%
\AgdaSymbol{(}\AgdaFunction{◇}\AgdaSymbol{)}\AgdaSpace{}%
\AgdaSymbol{((}\AgdaFunction{Forall}\AgdaSpace{}%
\AgdaSymbol{(}\AgdaBound{P}\AgdaSpace{}%
\AgdaOperator{\AgdaFunction{∧}}\AgdaSpace{}%
\AgdaBound{Q}\AgdaSymbol{))}\AgdaSpace{}%
\AgdaOperator{\AgdaFunction{⊃}}\AgdaSpace{}%
\AgdaSymbol{(}\AgdaFunction{Forall}\AgdaSpace{}%
\AgdaBound{P}\AgdaSpace{}%
\AgdaOperator{\AgdaFunction{∧}}\AgdaSpace{}%
\AgdaFunction{Forall}\AgdaSpace{}%
\AgdaBound{Q}\AgdaSymbol{))}\<%
\\
%
\>[2]\AgdaFunction{∀P∧Q⊃∀P∧∀Q}\AgdaSpace{}%
\AgdaBound{γ}\AgdaSpace{}%
\AgdaBound{x}\AgdaSpace{}%
\AgdaSymbol{=}\AgdaSpace{}%
\AgdaSymbol{(λ}\AgdaSpace{}%
\AgdaBound{d}\AgdaSpace{}%
\AgdaSymbol{→}\AgdaSpace{}%
\AgdaField{fst}\AgdaSpace{}%
\AgdaSymbol{(}\AgdaBound{x}\AgdaSpace{}%
\AgdaBound{d}\AgdaSymbol{))}\AgdaSpace{}%
\AgdaOperator{\AgdaInductiveConstructor{,p}}\AgdaSpace{}%
\AgdaSymbol{λ}\AgdaSpace{}%
\AgdaBound{d}\AgdaSpace{}%
\AgdaSymbol{→}\AgdaSpace{}%
\AgdaField{snd}\AgdaSpace{}%
\AgdaSymbol{(}\AgdaBound{x}\AgdaSpace{}%
\AgdaBound{d}\AgdaSymbol{)}\<%
\end{code}}
\begin{code}[hide]%
%
\>[2]\AgdaFunction{∀P∧∀Q⊃∀P∧Q}\AgdaSpace{}%
\AgdaSymbol{:}\AgdaSpace{}%
\AgdaSymbol{\{}\AgdaBound{P}\AgdaSpace{}%
\AgdaBound{Q}\AgdaSpace{}%
\AgdaSymbol{:}\AgdaSpace{}%
\AgdaFunction{For}\AgdaSpace{}%
\AgdaSymbol{(}\AgdaFunction{◇}\AgdaSpace{}%
\AgdaOperator{\AgdaFunction{▹ₜ}}\AgdaSymbol{)\}}\AgdaSpace{}%
\AgdaSymbol{→}\AgdaSpace{}%
\AgdaFunction{Pf}\AgdaSpace{}%
\AgdaSymbol{(}\AgdaFunction{◇}\AgdaSymbol{)}\AgdaSpace{}%
\AgdaSymbol{((}\AgdaFunction{Forall}\AgdaSpace{}%
\AgdaBound{P}\AgdaSpace{}%
\AgdaOperator{\AgdaFunction{∧}}\AgdaSpace{}%
\AgdaFunction{Forall}\AgdaSpace{}%
\AgdaBound{Q}\AgdaSymbol{)}\AgdaSpace{}%
\AgdaOperator{\AgdaFunction{⊃}}\AgdaSpace{}%
\AgdaSymbol{(}\AgdaFunction{Forall}\AgdaSpace{}%
\AgdaSymbol{(}\AgdaBound{P}\AgdaSpace{}%
\AgdaOperator{\AgdaFunction{∧}}\AgdaSpace{}%
\AgdaBound{Q}\AgdaSymbol{)))}\<%
\\
%
\>[2]\AgdaFunction{∀P∧∀Q⊃∀P∧Q}\AgdaSpace{}%
\AgdaSymbol{=}\AgdaSpace{}%
\AgdaSymbol{λ}\AgdaSpace{}%
\AgdaBound{γ}\AgdaSpace{}%
\AgdaBound{x}\AgdaSpace{}%
\AgdaBound{d}\AgdaSpace{}%
\AgdaSymbol{→}\AgdaSpace{}%
\AgdaField{fst}\AgdaSpace{}%
\AgdaBound{x}\AgdaSpace{}%
\AgdaBound{d}\AgdaSpace{}%
\AgdaOperator{\AgdaInductiveConstructor{,p}}\AgdaSpace{}%
\AgdaField{snd}\AgdaSpace{}%
\AgdaBound{x}\AgdaSpace{}%
\AgdaBound{d}\<%
\\
%
\\[\AgdaEmptyExtraSkip]%
\>[0]\AgdaKeyword{module}\AgdaSpace{}%
\AgdaModule{workInBool}\AgdaSpace{}%
\AgdaKeyword{where}\<%
\\
\>[0][@{}l@{\AgdaIndent{0}}]%
\>[2]\AgdaKeyword{import}\AgdaSpace{}%
\AgdaModule{BoolModel}\AgdaSpace{}%
\AgdaDatatype{funar}\AgdaSpace{}%
\AgdaDatatype{relar}\AgdaSpace{}%
\AgdaDatatype{ℕ}\AgdaSpace{}%
\AgdaSymbol{(λ}\AgdaSpace{}%
\AgdaBound{n}\AgdaSpace{}%
\AgdaBound{\AgdaUnderscore{}}\AgdaSpace{}%
\AgdaBound{\AgdaUnderscore{}}\AgdaSpace{}%
\AgdaSymbol{→}\AgdaSpace{}%
\AgdaBound{n}\AgdaSymbol{)}\AgdaSpace{}%
\AgdaSymbol{(λ}\AgdaSpace{}%
\AgdaBound{n}\AgdaSpace{}%
\AgdaBound{\AgdaUnderscore{}}\AgdaSpace{}%
\AgdaBound{\AgdaUnderscore{}}\AgdaSpace{}%
\AgdaSymbol{→}\AgdaSpace{}%
\AgdaInductiveConstructor{true}\AgdaSymbol{)}\ as\ \AgdaModule{B}\<%
\\
%
\>[2]\AgdaKeyword{open}\AgdaSpace{}%
\AgdaModule{Model}\AgdaSpace{}%
\AgdaFunction{B.B}\<%
\\
\>[0]\<%
\\
%
\>[2]\AgdaFunction{lem}\AgdaSpace{}%
\AgdaSymbol{:}\AgdaSpace{}%
\AgdaSymbol{∀\{}\AgdaBound{Γ}\AgdaSymbol{\}\{}\AgdaBound{A}\AgdaSpace{}%
\AgdaSymbol{:}\AgdaSpace{}%
\AgdaFunction{For}\AgdaSpace{}%
\AgdaBound{Γ}\AgdaSymbol{\}}\AgdaSpace{}%
\AgdaSymbol{→}\AgdaSpace{}%
\AgdaFunction{Pf}\AgdaSpace{}%
\AgdaBound{Γ}\AgdaSpace{}%
\AgdaSymbol{(}\AgdaBound{A}\AgdaSpace{}%
\AgdaOperator{\AgdaFunction{∨}}\AgdaSpace{}%
\AgdaSymbol{(}\AgdaBound{A}\AgdaSpace{}%
\AgdaOperator{\AgdaFunction{⊃}}\AgdaSpace{}%
\AgdaFunction{⊥}\AgdaSymbol{))}\<%
\\
%
\>[2]\AgdaFunction{lem}\AgdaSpace{}%
\AgdaSymbol{\{}\AgdaBound{Γ}\AgdaSymbol{\}}\AgdaSpace{}%
\AgdaSymbol{\{}\AgdaBound{A}\AgdaSymbol{\}}\AgdaSpace{}%
\AgdaBound{γ}\AgdaSpace{}%
\AgdaKeyword{with}\AgdaSpace{}%
\AgdaBound{A}\AgdaSpace{}%
\AgdaBound{γ}\<%
\\
%
\>[2]\AgdaSymbol{...}\AgdaSpace{}%
\AgdaSymbol{|}\AgdaSpace{}%
\AgdaInductiveConstructor{false}\AgdaSpace{}%
\AgdaSymbol{=}\AgdaSpace{}%
\AgdaInductiveConstructor{refl}\<%
\\
%
\>[2]\AgdaSymbol{...}\AgdaSpace{}%
\AgdaSymbol{|}\AgdaSpace{}%
\AgdaInductiveConstructor{true}\AgdaSpace{}%
\AgdaSymbol{=}\AgdaSpace{}%
\AgdaInductiveConstructor{refl}\<%
\end{code}
\newcommand{\pBooE}{
\begin{code}%
%
\>[2]\AgdaFunction{A⊃B⊃A}\AgdaSpace{}%
\AgdaSymbol{:}\AgdaSpace{}%
\AgdaSymbol{\{}\AgdaBound{A}\AgdaSpace{}%
\AgdaBound{B}\AgdaSpace{}%
\AgdaSymbol{:}\AgdaSpace{}%
\AgdaFunction{For}\AgdaSpace{}%
\AgdaFunction{◇}\AgdaSymbol{\}}\AgdaSpace{}%
\AgdaSymbol{→}\AgdaSpace{}%
\AgdaFunction{Pf}\AgdaSpace{}%
\AgdaSymbol{(}\AgdaFunction{◇}\AgdaSymbol{)}\AgdaSpace{}%
\AgdaSymbol{(}\AgdaBound{A}\AgdaSpace{}%
\AgdaOperator{\AgdaFunction{⊃}}\AgdaSpace{}%
\AgdaBound{B}\AgdaSpace{}%
\AgdaOperator{\AgdaFunction{⊃}}\AgdaSpace{}%
\AgdaBound{A}\AgdaSymbol{)}\<%
\\
%
\>[2]\AgdaFunction{A⊃B⊃A}\AgdaSpace{}%
\AgdaSymbol{\{}\AgdaBound{A}\AgdaSymbol{\}}\AgdaSpace{}%
\AgdaSymbol{\{}\AgdaBound{B}\AgdaSymbol{\}}\AgdaSpace{}%
\AgdaBound{γ}\AgdaSpace{}%
\AgdaKeyword{with}\AgdaSpace{}%
\AgdaBound{A}\AgdaSpace{}%
\AgdaBound{γ}\AgdaSpace{}%
\AgdaSymbol{|}\AgdaSpace{}%
\AgdaBound{B}\AgdaSpace{}%
\AgdaBound{γ}\<%
\\
%
\>[2]\AgdaSymbol{...}\AgdaSpace{}%
\AgdaSymbol{|}\AgdaSpace{}%
\AgdaInductiveConstructor{false}\AgdaSpace{}%
\AgdaSymbol{|}\AgdaSpace{}%
\AgdaInductiveConstructor{false}\AgdaSpace{}%
\AgdaSymbol{=}\AgdaSpace{}%
\AgdaInductiveConstructor{refl}\<%
\\
%
\>[2]\AgdaSymbol{...}\AgdaSpace{}%
\AgdaSymbol{|}\AgdaSpace{}%
\AgdaInductiveConstructor{false}\AgdaSpace{}%
\AgdaSymbol{|}\AgdaSpace{}%
\AgdaInductiveConstructor{true}\AgdaSpace{}%
\AgdaSymbol{=}\AgdaSpace{}%
\AgdaInductiveConstructor{refl}\<%
\\
%
\>[2]\AgdaSymbol{...}\AgdaSpace{}%
\AgdaSymbol{|}\AgdaSpace{}%
\AgdaInductiveConstructor{true}\AgdaSpace{}%
\AgdaSymbol{|}\AgdaSpace{}%
\AgdaInductiveConstructor{false}\AgdaSpace{}%
\AgdaSymbol{=}\AgdaSpace{}%
\AgdaInductiveConstructor{refl}\<%
\\
%
\>[2]\AgdaSymbol{...}\AgdaSpace{}%
\AgdaSymbol{|}\AgdaSpace{}%
\AgdaInductiveConstructor{true}\AgdaSpace{}%
\AgdaSymbol{|}\AgdaSpace{}%
\AgdaInductiveConstructor{true}\AgdaSpace{}%
\AgdaSymbol{=}\AgdaSpace{}%
\AgdaInductiveConstructor{refl}\<%
\end{code}}
\begin{code}[hide]%
%
\>[2]\AgdaFunction{t\&\&t→t₁}\AgdaSpace{}%
\AgdaSymbol{:}\AgdaSpace{}%
\AgdaSymbol{∀\{}\AgdaBound{a}\AgdaSpace{}%
\AgdaBound{b}\AgdaSymbol{\}}\AgdaSpace{}%
\AgdaSymbol{→}\AgdaSpace{}%
\AgdaBound{a}\AgdaSpace{}%
\AgdaOperator{\AgdaFunction{\&\&}}\AgdaSpace{}%
\AgdaBound{b}\AgdaSpace{}%
\AgdaOperator{\AgdaDatatype{≡}}\AgdaSpace{}%
\AgdaInductiveConstructor{true}\AgdaSpace{}%
\AgdaSymbol{→}\AgdaSpace{}%
\AgdaBound{a}\AgdaSpace{}%
\AgdaOperator{\AgdaDatatype{≡}}\AgdaSpace{}%
\AgdaInductiveConstructor{true}\<%
\\
%
\>[2]\AgdaFunction{t\&\&t→t₁}\AgdaSpace{}%
\AgdaSymbol{\{}\AgdaInductiveConstructor{true}\AgdaSymbol{\}}\AgdaSpace{}%
\AgdaSymbol{\{}\AgdaInductiveConstructor{true}\AgdaSymbol{\}}\AgdaSpace{}%
\AgdaInductiveConstructor{refl}\AgdaSpace{}%
\AgdaSymbol{=}\AgdaSpace{}%
\AgdaInductiveConstructor{refl}\<%
\\
%
\>[2]\AgdaFunction{t\&\&t→t₁}\AgdaSpace{}%
\AgdaSymbol{\{}\AgdaInductiveConstructor{false}\AgdaSymbol{\}}\AgdaSpace{}%
\AgdaSymbol{\{}\AgdaInductiveConstructor{true}\AgdaSymbol{\}}\AgdaSpace{}%
\AgdaBound{x}\AgdaSpace{}%
\AgdaSymbol{=}\AgdaSpace{}%
\AgdaBound{x}\<%
\\
%
\>[2]\AgdaFunction{t\&\&t→t₂}\AgdaSpace{}%
\AgdaSymbol{:}\AgdaSpace{}%
\AgdaSymbol{∀\{}\AgdaBound{a}\AgdaSpace{}%
\AgdaBound{b}\AgdaSymbol{\}}\AgdaSpace{}%
\AgdaSymbol{→}\AgdaSpace{}%
\AgdaBound{a}\AgdaSpace{}%
\AgdaOperator{\AgdaFunction{\&\&}}\AgdaSpace{}%
\AgdaBound{b}\AgdaSpace{}%
\AgdaOperator{\AgdaDatatype{≡}}\AgdaSpace{}%
\AgdaInductiveConstructor{true}\AgdaSpace{}%
\AgdaSymbol{→}\AgdaSpace{}%
\AgdaBound{b}\AgdaSpace{}%
\AgdaOperator{\AgdaDatatype{≡}}\AgdaSpace{}%
\AgdaInductiveConstructor{true}\<%
\\
%
\>[2]\AgdaFunction{t\&\&t→t₂}\AgdaSpace{}%
\AgdaSymbol{\{}\AgdaInductiveConstructor{true}\AgdaSymbol{\}}\AgdaSpace{}%
\AgdaSymbol{\{}\AgdaInductiveConstructor{true}\AgdaSymbol{\}}\AgdaSpace{}%
\AgdaInductiveConstructor{refl}\AgdaSpace{}%
\AgdaSymbol{=}\AgdaSpace{}%
\AgdaInductiveConstructor{refl}\<%
\\
%
\>[2]\AgdaFunction{t\&\&t→t₂}\AgdaSpace{}%
\AgdaSymbol{\{}\AgdaInductiveConstructor{true}\AgdaSymbol{\}}\AgdaSpace{}%
\AgdaSymbol{\{}\AgdaInductiveConstructor{false}\AgdaSymbol{\}}\AgdaSpace{}%
\AgdaBound{x}\AgdaSpace{}%
\AgdaSymbol{=}\AgdaSpace{}%
\AgdaBound{x}\<%
\end{code}
\newcommand{\pBooC}{
\begin{code}%
%
\>[2]\AgdaFunction{∀P∧Q⊃∀P∧∀Q}\AgdaSpace{}%
\AgdaSymbol{:}%
\>[1109I]\AgdaSymbol{\{}\AgdaBound{P}\AgdaSpace{}%
\AgdaBound{Q}\AgdaSpace{}%
\AgdaSymbol{:}\AgdaSpace{}%
\AgdaFunction{For}\AgdaSpace{}%
\AgdaSymbol{(}\AgdaFunction{◇}\AgdaSpace{}%
\AgdaOperator{\AgdaFunction{▹ₜ}}\AgdaSymbol{)\}}\AgdaSpace{}%
\AgdaSymbol{→}\<%
\\
\>[1109I][@{}l@{\AgdaIndent{0}}]%
\>[16]\AgdaFunction{Pf}\AgdaSpace{}%
\AgdaSymbol{(}\AgdaFunction{◇}\AgdaSymbol{)}\AgdaSpace{}%
\AgdaSymbol{((}\AgdaFunction{Forall}\AgdaSpace{}%
\AgdaSymbol{(}\AgdaBound{P}\AgdaSpace{}%
\AgdaOperator{\AgdaFunction{∧}}\AgdaSpace{}%
\AgdaBound{Q}\AgdaSymbol{))}\AgdaSpace{}%
\AgdaOperator{\AgdaFunction{⊃}}\AgdaSpace{}%
\AgdaSymbol{(}\AgdaFunction{Forall}\AgdaSpace{}%
\AgdaBound{P}\AgdaSpace{}%
\AgdaOperator{\AgdaFunction{∧}}\AgdaSpace{}%
\AgdaFunction{Forall}\AgdaSpace{}%
\AgdaBound{Q}\AgdaSymbol{))}\<%
\\
%
\>[2]\AgdaFunction{∀P∧Q⊃∀P∧∀Q}\AgdaSpace{}%
\AgdaSymbol{\{}\AgdaBound{P}\AgdaSymbol{\}}\AgdaSpace{}%
\AgdaSymbol{\{}\AgdaBound{Q}\AgdaSymbol{\}}\AgdaSpace{}%
\AgdaBound{γ}\AgdaSpace{}%
\AgdaSymbol{=}\<%
\\
\>[2][@{}l@{\AgdaIndent{0}}]%
\>[4]\AgdaFunction{ind⊎p}\<%
\\
\>[4][@{}l@{\AgdaIndent{0}}]%
\>[6]\AgdaSymbol{(λ}\AgdaSpace{}%
\AgdaBound{x}\AgdaSpace{}%
\AgdaSymbol{→}\AgdaSpace{}%
\AgdaFunction{case}\AgdaSpace{}%
\AgdaSymbol{(λ}\AgdaSpace{}%
\AgdaBound{\AgdaUnderscore{}}\AgdaSpace{}%
\AgdaSymbol{→}\AgdaSpace{}%
\AgdaInductiveConstructor{true}\AgdaSymbol{)}\AgdaSpace{}%
\AgdaSymbol{(λ}\AgdaSpace{}%
\AgdaBound{\AgdaUnderscore{}}\AgdaSpace{}%
\AgdaSymbol{→}\AgdaSpace{}%
\AgdaInductiveConstructor{false}\AgdaSymbol{)}\<%
\\
\>[6][@{}l@{\AgdaIndent{0}}]%
\>[8]\AgdaBound{x}\AgdaSpace{}%
\AgdaOperator{\AgdaFunction{=>}}\AgdaSpace{}%
\AgdaSymbol{(}\AgdaFunction{Forall}\AgdaSpace{}%
\AgdaBound{P}\AgdaSpace{}%
\AgdaOperator{\AgdaFunction{∧}}\AgdaSpace{}%
\AgdaFunction{Forall}\AgdaSpace{}%
\AgdaBound{Q}\AgdaSymbol{)}\AgdaSpace{}%
\AgdaBound{γ}\AgdaSpace{}%
\AgdaOperator{\AgdaDatatype{≡}}\AgdaSpace{}%
\AgdaInductiveConstructor{true}\AgdaSymbol{)}\<%
\\
%
\>[6]\AgdaSymbol{(λ}\AgdaSpace{}%
\AgdaBound{h₁}\AgdaSpace{}%
\AgdaSymbol{→}\AgdaSpace{}%
\AgdaFunction{ind⊎p}\<%
\\
\>[6][@{}l@{\AgdaIndent{0}}]%
\>[8]\AgdaSymbol{(λ}%
\>[1154I]\AgdaBound{x}\AgdaSpace{}%
\AgdaSymbol{→}\AgdaSpace{}%
\AgdaInductiveConstructor{true}\AgdaSpace{}%
\AgdaOperator{\AgdaFunction{=>}}\AgdaSpace{}%
\AgdaSymbol{(}\AgdaFunction{case}\AgdaSpace{}%
\AgdaSymbol{(λ}\AgdaSpace{}%
\AgdaBound{\AgdaUnderscore{}}\AgdaSpace{}%
\AgdaSymbol{→}\AgdaSpace{}%
\AgdaInductiveConstructor{true}\AgdaSymbol{)}\AgdaSpace{}%
\AgdaSymbol{(λ}\AgdaSpace{}%
\AgdaBound{\AgdaUnderscore{}}\AgdaSpace{}%
\AgdaSymbol{→}\AgdaSpace{}%
\AgdaInductiveConstructor{false}\AgdaSymbol{)}\<%
\\
\>[1154I][@{}l@{\AgdaIndent{0}}]%
\>[12]\AgdaBound{x}\AgdaSpace{}%
\AgdaOperator{\AgdaFunction{\&\&}}\AgdaSpace{}%
\AgdaFunction{Forall}\AgdaSpace{}%
\AgdaBound{Q}\AgdaSpace{}%
\AgdaBound{γ}\AgdaSymbol{)}\<%
\\
\>[8][@{}l@{\AgdaIndent{0}}]%
\>[10]\AgdaOperator{\AgdaDatatype{≡}}\AgdaSpace{}%
\AgdaInductiveConstructor{true}\AgdaSymbol{)}\<%
\\
%
\>[8]\AgdaSymbol{(λ}\AgdaSpace{}%
\AgdaBound{h₂}\AgdaSpace{}%
\AgdaSymbol{→}\AgdaSpace{}%
\AgdaFunction{ind⊎p}\<%
\\
\>[8][@{}l@{\AgdaIndent{0}}]%
\>[10]\AgdaSymbol{(λ}\AgdaSpace{}%
\AgdaBound{x}\AgdaSpace{}%
\AgdaSymbol{→}\AgdaSpace{}%
\AgdaInductiveConstructor{true}\AgdaSpace{}%
\AgdaOperator{\AgdaFunction{=>}}\AgdaSpace{}%
\AgdaSymbol{(}\AgdaInductiveConstructor{true}\AgdaSpace{}%
\AgdaOperator{\AgdaFunction{\&\&}}\AgdaSpace{}%
\AgdaFunction{case}\AgdaSpace{}%
\AgdaSymbol{(λ}\AgdaSpace{}%
\AgdaBound{\AgdaUnderscore{}}\AgdaSpace{}%
\AgdaSymbol{→}\AgdaSpace{}%
\AgdaInductiveConstructor{true}\AgdaSymbol{)}\AgdaSpace{}%
\AgdaSymbol{(λ}\AgdaSpace{}%
\AgdaBound{\AgdaUnderscore{}}\AgdaSpace{}%
\AgdaSymbol{→}\AgdaSpace{}%
\AgdaInductiveConstructor{false}\AgdaSymbol{)}\AgdaSpace{}%
\AgdaBound{x}\AgdaSymbol{)}\<%
\\
\>[10][@{}l@{\AgdaIndent{0}}]%
\>[12]\AgdaOperator{\AgdaDatatype{≡}}\AgdaSpace{}%
\AgdaInductiveConstructor{true}\AgdaSymbol{)}\<%
\\
%
\>[10]\AgdaSymbol{(λ}\AgdaSpace{}%
\AgdaBound{h₃}\AgdaSpace{}%
\AgdaSymbol{→}\AgdaSpace{}%
\AgdaInductiveConstructor{refl}\AgdaSymbol{)}\<%
\\
%
\>[10]\AgdaSymbol{(λ}\AgdaSpace{}%
\AgdaBound{h₃}\AgdaSpace{}%
\AgdaSymbol{→}\AgdaSpace{}%
\AgdaFunction{exfalsop}\AgdaSpace{}%
\AgdaSymbol{(}\AgdaBound{h₃}\AgdaSpace{}%
\AgdaSymbol{λ}\AgdaSpace{}%
\AgdaBound{d}\AgdaSpace{}%
\AgdaSymbol{→}\AgdaSpace{}%
\AgdaFunction{t\&\&t→t₂}\AgdaSpace{}%
\AgdaSymbol{(}\AgdaBound{h₁}\AgdaSpace{}%
\AgdaBound{d}\AgdaSymbol{)))}\<%
\\
%
\>[10]\AgdaSymbol{(}\AgdaFunction{B.lem}\AgdaSpace{}%
\AgdaSymbol{((}\AgdaBound{d}\AgdaSpace{}%
\AgdaSymbol{:}\AgdaSpace{}%
\AgdaDatatype{ℕ}\AgdaSymbol{)}\AgdaSpace{}%
\AgdaSymbol{→}\AgdaSpace{}%
\AgdaBound{Q}\AgdaSpace{}%
\AgdaSymbol{(}\AgdaBound{γ}\AgdaSpace{}%
\AgdaOperator{\AgdaInductiveConstructor{,}}\AgdaSpace{}%
\AgdaBound{d}\AgdaSymbol{)}\AgdaSpace{}%
\AgdaOperator{\AgdaDatatype{≡}}\AgdaSpace{}%
\AgdaInductiveConstructor{true}\AgdaSymbol{)))}\<%
\\
%
\>[8]\AgdaSymbol{(λ}\AgdaSpace{}%
\AgdaBound{h₂}\AgdaSpace{}%
\AgdaSymbol{→}\AgdaSpace{}%
\AgdaFunction{ind⊎p}\<%
\\
\>[8][@{}l@{\AgdaIndent{0}}]%
\>[10]\AgdaSymbol{(λ}%
\>[1218I]\AgdaBound{x}\AgdaSpace{}%
\AgdaSymbol{→}\AgdaSpace{}%
\AgdaInductiveConstructor{true}\AgdaSpace{}%
\AgdaOperator{\AgdaFunction{=>}}\AgdaSpace{}%
\AgdaSymbol{(}\AgdaInductiveConstructor{false}\AgdaSpace{}%
\AgdaOperator{\AgdaFunction{\&\&}}\AgdaSpace{}%
\AgdaFunction{case}\AgdaSpace{}%
\AgdaSymbol{(λ}\AgdaSpace{}%
\AgdaBound{\AgdaUnderscore{}}\AgdaSpace{}%
\AgdaSymbol{→}\AgdaSpace{}%
\AgdaInductiveConstructor{true}\AgdaSymbol{)}\AgdaSpace{}%
\AgdaSymbol{(λ}\AgdaSpace{}%
\AgdaBound{\AgdaUnderscore{}}\AgdaSpace{}%
\AgdaSymbol{→}\AgdaSpace{}%
\AgdaInductiveConstructor{false}\AgdaSymbol{)}\<%
\\
\>[1218I][@{}l@{\AgdaIndent{0}}]%
\>[14]\AgdaBound{x}\AgdaSymbol{)}\<%
\\
\>[10][@{}l@{\AgdaIndent{0}}]%
\>[12]\AgdaOperator{\AgdaDatatype{≡}}\AgdaSpace{}%
\AgdaInductiveConstructor{true}\AgdaSymbol{)}\<%
\\
%
\>[10]\AgdaSymbol{(λ}\AgdaSpace{}%
\AgdaBound{h₃}\AgdaSpace{}%
\AgdaSymbol{→}\AgdaSpace{}%
\AgdaFunction{exfalsop}\AgdaSpace{}%
\AgdaSymbol{(}\AgdaBound{h₂}\AgdaSpace{}%
\AgdaSymbol{λ}\AgdaSpace{}%
\AgdaBound{d}\AgdaSpace{}%
\AgdaSymbol{→}\AgdaSpace{}%
\AgdaFunction{t\&\&t→t₁}\AgdaSpace{}%
\AgdaSymbol{(}\AgdaBound{h₁}\AgdaSpace{}%
\AgdaBound{d}\AgdaSymbol{)))}\<%
\\
%
\>[10]\AgdaSymbol{(λ}\AgdaSpace{}%
\AgdaBound{h₃}\AgdaSpace{}%
\AgdaSymbol{→}\AgdaSpace{}%
\AgdaFunction{exfalsop}\AgdaSpace{}%
\AgdaSymbol{(}\AgdaBound{h₃}\AgdaSpace{}%
\AgdaSymbol{λ}\AgdaSpace{}%
\AgdaBound{d}\AgdaSpace{}%
\AgdaSymbol{→}\AgdaSpace{}%
\AgdaFunction{t\&\&t→t₂}\AgdaSpace{}%
\AgdaSymbol{(}\AgdaBound{h₁}\AgdaSpace{}%
\AgdaBound{d}\AgdaSymbol{)))}\<%
\\
%
\>[10]\AgdaSymbol{(}\AgdaFunction{B.lem}\AgdaSpace{}%
\AgdaSymbol{((}\AgdaBound{d}\AgdaSpace{}%
\AgdaSymbol{:}\AgdaSpace{}%
\AgdaDatatype{ℕ}\AgdaSymbol{)}\AgdaSpace{}%
\AgdaSymbol{→}\AgdaSpace{}%
\AgdaBound{Q}\AgdaSpace{}%
\AgdaSymbol{(}\AgdaBound{γ}\AgdaSpace{}%
\AgdaOperator{\AgdaInductiveConstructor{,}}\AgdaSpace{}%
\AgdaBound{d}\AgdaSymbol{)}\AgdaSpace{}%
\AgdaOperator{\AgdaDatatype{≡}}\AgdaSpace{}%
\AgdaInductiveConstructor{true}\AgdaSymbol{)))}\<%
\\
%
\>[8]\AgdaSymbol{(}\AgdaFunction{B.lem}\AgdaSpace{}%
\AgdaSymbol{((}\AgdaBound{d}\AgdaSpace{}%
\AgdaSymbol{:}\AgdaSpace{}%
\AgdaDatatype{ℕ}\AgdaSymbol{)}\AgdaSpace{}%
\AgdaSymbol{→}\AgdaSpace{}%
\AgdaBound{P}\AgdaSpace{}%
\AgdaSymbol{(}\AgdaBound{γ}\AgdaSpace{}%
\AgdaOperator{\AgdaInductiveConstructor{,}}\AgdaSpace{}%
\AgdaBound{d}\AgdaSymbol{)}\AgdaSpace{}%
\AgdaOperator{\AgdaDatatype{≡}}\AgdaSpace{}%
\AgdaInductiveConstructor{true}\AgdaSymbol{)))}\<%
\\
%
\>[6]\AgdaSymbol{(λ}\AgdaSpace{}%
\AgdaBound{h₁}\AgdaSpace{}%
\AgdaSymbol{→}\AgdaSpace{}%
\AgdaFunction{ind⊎p}\<%
\\
\>[6][@{}l@{\AgdaIndent{0}}]%
\>[8]\AgdaSymbol{(λ}%
\>[1277I]\AgdaBound{x}\AgdaSpace{}%
\AgdaSymbol{→}\AgdaSpace{}%
\AgdaInductiveConstructor{false}\AgdaSpace{}%
\AgdaOperator{\AgdaFunction{=>}}\AgdaSpace{}%
\AgdaSymbol{(}\AgdaFunction{case}\AgdaSpace{}%
\AgdaSymbol{(λ}\AgdaSpace{}%
\AgdaBound{\AgdaUnderscore{}}\AgdaSpace{}%
\AgdaSymbol{→}\AgdaSpace{}%
\AgdaInductiveConstructor{true}\AgdaSymbol{)}\AgdaSpace{}%
\AgdaSymbol{(λ}\AgdaSpace{}%
\AgdaBound{\AgdaUnderscore{}}\AgdaSpace{}%
\AgdaSymbol{→}\AgdaSpace{}%
\AgdaInductiveConstructor{false}\AgdaSymbol{)}\<%
\\
\>[1277I][@{}l@{\AgdaIndent{0}}]%
\>[12]\AgdaBound{x}\AgdaSpace{}%
\AgdaOperator{\AgdaFunction{\&\&}}\AgdaSpace{}%
\AgdaFunction{Forall}\AgdaSpace{}%
\AgdaBound{Q}\AgdaSpace{}%
\AgdaBound{γ}\AgdaSymbol{)}\<%
\\
\>[8][@{}l@{\AgdaIndent{0}}]%
\>[10]\AgdaOperator{\AgdaDatatype{≡}}\AgdaSpace{}%
\AgdaInductiveConstructor{true}\AgdaSymbol{)}\<%
\\
%
\>[8]\AgdaSymbol{(λ}\AgdaSpace{}%
\AgdaBound{h₂}\AgdaSpace{}%
\AgdaSymbol{→}\AgdaSpace{}%
\AgdaFunction{ind⊎p}\<%
\\
\>[8][@{}l@{\AgdaIndent{0}}]%
\>[10]\AgdaSymbol{(λ}\AgdaSpace{}%
\AgdaBound{x}\AgdaSpace{}%
\AgdaSymbol{→}\AgdaSpace{}%
\AgdaInductiveConstructor{false}\AgdaSpace{}%
\AgdaOperator{\AgdaFunction{=>}}\AgdaSpace{}%
\AgdaSymbol{(}\AgdaInductiveConstructor{true}\AgdaSpace{}%
\AgdaOperator{\AgdaFunction{\&\&}}\AgdaSpace{}%
\AgdaFunction{case}\AgdaSpace{}%
\AgdaSymbol{(λ}\AgdaSpace{}%
\AgdaBound{\AgdaUnderscore{}}\AgdaSpace{}%
\AgdaSymbol{→}\AgdaSpace{}%
\AgdaInductiveConstructor{true}\AgdaSymbol{)}\AgdaSpace{}%
\AgdaSymbol{(λ}\AgdaSpace{}%
\AgdaBound{\AgdaUnderscore{}}\AgdaSpace{}%
\AgdaSymbol{→}\AgdaSpace{}%
\AgdaInductiveConstructor{false}\AgdaSymbol{)}\<%
\\
\>[10][@{}l@{\AgdaIndent{0}}]%
\>[12]\AgdaBound{x}\AgdaSymbol{)}\AgdaSpace{}%
\AgdaOperator{\AgdaDatatype{≡}}\AgdaSpace{}%
\AgdaInductiveConstructor{true}\AgdaSymbol{)}\<%
\\
%
\>[10]\AgdaSymbol{(λ}\AgdaSpace{}%
\AgdaBound{h₃}\AgdaSpace{}%
\AgdaSymbol{→}\AgdaSpace{}%
\AgdaInductiveConstructor{refl}\AgdaSymbol{)}\<%
\\
%
\>[10]\AgdaSymbol{(λ}\AgdaSpace{}%
\AgdaBound{h₃}\AgdaSpace{}%
\AgdaSymbol{→}\AgdaSpace{}%
\AgdaInductiveConstructor{refl}\AgdaSymbol{)}\<%
\\
%
\>[10]\AgdaSymbol{(}\AgdaFunction{B.lem}\AgdaSpace{}%
\AgdaSymbol{((}\AgdaBound{d}\AgdaSpace{}%
\AgdaSymbol{:}\AgdaSpace{}%
\AgdaDatatype{ℕ}\AgdaSymbol{)}\AgdaSpace{}%
\AgdaSymbol{→}\AgdaSpace{}%
\AgdaBound{Q}\AgdaSpace{}%
\AgdaSymbol{(}\AgdaBound{γ}\AgdaSpace{}%
\AgdaOperator{\AgdaInductiveConstructor{,}}\AgdaSpace{}%
\AgdaBound{d}\AgdaSymbol{)}\AgdaSpace{}%
\AgdaOperator{\AgdaDatatype{≡}}\AgdaSpace{}%
\AgdaInductiveConstructor{true}\AgdaSymbol{)))}\<%
\\
%
\>[8]\AgdaSymbol{(λ}\AgdaSpace{}%
\AgdaBound{h₂}\AgdaSpace{}%
\AgdaSymbol{→}\AgdaSpace{}%
\AgdaFunction{ind⊎p}\<%
\\
\>[8][@{}l@{\AgdaIndent{0}}]%
\>[10]\AgdaSymbol{(λ}\AgdaSpace{}%
\AgdaBound{x}\AgdaSpace{}%
\AgdaSymbol{→}\AgdaSpace{}%
\AgdaInductiveConstructor{false}\AgdaSpace{}%
\AgdaOperator{\AgdaFunction{=>}}\AgdaSpace{}%
\AgdaSymbol{(}\AgdaInductiveConstructor{false}\AgdaSpace{}%
\AgdaOperator{\AgdaFunction{\&\&}}\AgdaSpace{}%
\AgdaFunction{case}\AgdaSpace{}%
\AgdaSymbol{(λ}\AgdaSpace{}%
\AgdaBound{\AgdaUnderscore{}}\AgdaSpace{}%
\AgdaSymbol{→}\AgdaSpace{}%
\AgdaInductiveConstructor{true}\AgdaSymbol{)}\AgdaSpace{}%
\AgdaSymbol{(λ}\AgdaSpace{}%
\AgdaBound{\AgdaUnderscore{}}\AgdaSpace{}%
\AgdaSymbol{→}\AgdaSpace{}%
\AgdaInductiveConstructor{false}\AgdaSymbol{)}\<%
\\
\>[10][@{}l@{\AgdaIndent{0}}]%
\>[12]\AgdaBound{x}\AgdaSymbol{)}\AgdaSpace{}%
\AgdaOperator{\AgdaDatatype{≡}}\AgdaSpace{}%
\AgdaInductiveConstructor{true}\AgdaSymbol{)}\<%
\\
%
\>[10]\AgdaSymbol{(λ}\AgdaSpace{}%
\AgdaBound{h₃}\AgdaSpace{}%
\AgdaSymbol{→}\AgdaSpace{}%
\AgdaInductiveConstructor{refl}\AgdaSymbol{)}\<%
\\
%
\>[10]\AgdaSymbol{(λ}\AgdaSpace{}%
\AgdaBound{h₃}\AgdaSpace{}%
\AgdaSymbol{→}\AgdaSpace{}%
\AgdaInductiveConstructor{refl}\AgdaSymbol{)}\<%
\\
%
\>[10]\AgdaSymbol{(}\AgdaFunction{B.lem}\AgdaSpace{}%
\AgdaSymbol{((}\AgdaBound{d}\AgdaSpace{}%
\AgdaSymbol{:}\AgdaSpace{}%
\AgdaDatatype{ℕ}\AgdaSymbol{)}\AgdaSpace{}%
\AgdaSymbol{→}\AgdaSpace{}%
\AgdaBound{Q}\AgdaSpace{}%
\AgdaSymbol{(}\AgdaBound{γ}\AgdaSpace{}%
\AgdaOperator{\AgdaInductiveConstructor{,}}\AgdaSpace{}%
\AgdaBound{d}\AgdaSymbol{)}\AgdaSpace{}%
\AgdaOperator{\AgdaDatatype{≡}}\AgdaSpace{}%
\AgdaInductiveConstructor{true}\AgdaSymbol{)))}\<%
\\
%
\>[8]\AgdaSymbol{(}\AgdaFunction{B.lem}\AgdaSpace{}%
\AgdaSymbol{((}\AgdaBound{d}\AgdaSpace{}%
\AgdaSymbol{:}\AgdaSpace{}%
\AgdaDatatype{ℕ}\AgdaSymbol{)}\AgdaSpace{}%
\AgdaSymbol{→}\AgdaSpace{}%
\AgdaBound{P}\AgdaSpace{}%
\AgdaSymbol{(}\AgdaBound{γ}\AgdaSpace{}%
\AgdaOperator{\AgdaInductiveConstructor{,}}\AgdaSpace{}%
\AgdaBound{d}\AgdaSymbol{)}\AgdaSpace{}%
\AgdaOperator{\AgdaDatatype{≡}}\AgdaSpace{}%
\AgdaInductiveConstructor{true}\AgdaSymbol{)))}\<%
\\
%
\>[6]\AgdaSymbol{(}\AgdaFunction{B.lem}\AgdaSpace{}%
\AgdaSymbol{((}\AgdaBound{d}\AgdaSpace{}%
\AgdaSymbol{:}\AgdaSpace{}%
\AgdaDatatype{ℕ}\AgdaSymbol{)}\AgdaSpace{}%
\AgdaSymbol{→}\AgdaSpace{}%
\AgdaSymbol{(}\AgdaBound{P}\AgdaSpace{}%
\AgdaSymbol{(}\AgdaBound{γ}\AgdaSpace{}%
\AgdaOperator{\AgdaInductiveConstructor{,}}\AgdaSpace{}%
\AgdaBound{d}\AgdaSymbol{)}\AgdaSpace{}%
\AgdaOperator{\AgdaFunction{\&\&}}\AgdaSpace{}%
\AgdaBound{Q}\AgdaSpace{}%
\AgdaSymbol{(}\AgdaBound{γ}\AgdaSpace{}%
\AgdaOperator{\AgdaInductiveConstructor{,}}\AgdaSpace{}%
\AgdaBound{d}\AgdaSymbol{))}\AgdaSpace{}%
\AgdaOperator{\AgdaDatatype{≡}}\AgdaSpace{}%
\AgdaInductiveConstructor{true}\AgdaSymbol{))}\<%
\end{code}}

\begin{code}[hide]%
\>[0]\AgdaKeyword{module}\AgdaSpace{}%
\AgdaModule{workInFamily}\AgdaSpace{}%
\AgdaKeyword{where}\<%
\\
\>[0][@{}l@{\AgdaIndent{0}}]%
\>[2]\AgdaFunction{ℕfun}\AgdaSpace{}%
\AgdaSymbol{:}\AgdaSpace{}%
\AgdaSymbol{(}\AgdaBound{n}\AgdaSpace{}%
\AgdaSymbol{:}\AgdaSpace{}%
\AgdaDatatype{ℕ}\AgdaSymbol{)}\AgdaSpace{}%
\AgdaSymbol{→}\AgdaSpace{}%
\AgdaDatatype{funar}\AgdaSpace{}%
\AgdaBound{n}\AgdaSpace{}%
\AgdaSymbol{→}\AgdaSpace{}%
\AgdaSymbol{(}\AgdaBound{i}\AgdaSpace{}%
\AgdaSymbol{:}\AgdaSpace{}%
\AgdaDatatype{ℕ}\AgdaSymbol{)}\AgdaSpace{}%
\AgdaSymbol{→}\AgdaSpace{}%
\AgdaDatatype{ℕ}\AgdaSpace{}%
\AgdaOperator{\AgdaFunction{\textasciicircum{}}}\AgdaSpace{}%
\AgdaBound{n}\AgdaSpace{}%
\AgdaSymbol{→}\AgdaSpace{}%
\AgdaDatatype{ℕ}\<%
\\
%
\>[2]\AgdaFunction{ℕfun}\AgdaSpace{}%
\AgdaSymbol{\AgdaUnderscore{}}\AgdaSpace{}%
\AgdaInductiveConstructor{zero}\AgdaSpace{}%
\AgdaSymbol{\AgdaUnderscore{}}\AgdaSpace{}%
\AgdaSymbol{\AgdaUnderscore{}}\AgdaSpace{}%
\AgdaSymbol{=}\AgdaSpace{}%
\AgdaInductiveConstructor{zero}\<%
\\
%
\>[2]\AgdaFunction{ℕfun}\AgdaSpace{}%
\AgdaSymbol{\AgdaUnderscore{}}\AgdaSpace{}%
\AgdaInductiveConstructor{suc}\AgdaSpace{}%
\AgdaSymbol{\AgdaUnderscore{}}\AgdaSpace{}%
\AgdaSymbol{(}\AgdaBound{n}\AgdaSpace{}%
\AgdaOperator{\AgdaInductiveConstructor{,}}\AgdaSpace{}%
\AgdaSymbol{\AgdaUnderscore{})}\AgdaSpace{}%
\AgdaSymbol{=}\AgdaSpace{}%
\AgdaInductiveConstructor{suc}\AgdaSpace{}%
\AgdaBound{n}\<%
\\
%
\>[2]\AgdaFunction{ℕfun}\AgdaSpace{}%
\AgdaSymbol{\AgdaUnderscore{}}\AgdaSpace{}%
\AgdaInductiveConstructor{+'}\AgdaSpace{}%
\AgdaSymbol{\AgdaUnderscore{}}\AgdaSpace{}%
\AgdaSymbol{(}\AgdaBound{m}\AgdaSpace{}%
\AgdaOperator{\AgdaInductiveConstructor{,}}\AgdaSpace{}%
\AgdaBound{n}\AgdaSpace{}%
\AgdaOperator{\AgdaInductiveConstructor{,}}\AgdaSpace{}%
\AgdaSymbol{\AgdaUnderscore{})}\AgdaSpace{}%
\AgdaSymbol{=}\AgdaSpace{}%
\AgdaBound{m}\AgdaSpace{}%
\AgdaOperator{\AgdaPrimitive{+}}\AgdaSpace{}%
\AgdaBound{n}\<%
\\
%
\\[\AgdaEmptyExtraSkip]%
%
\>[2]\AgdaFunction{ℕRel}\AgdaSpace{}%
\AgdaSymbol{:}\AgdaSpace{}%
\AgdaSymbol{(}\AgdaBound{n}\AgdaSpace{}%
\AgdaSymbol{:}\AgdaSpace{}%
\AgdaDatatype{ℕ}\AgdaSymbol{)}\AgdaSpace{}%
\AgdaSymbol{→}\AgdaSpace{}%
\AgdaDatatype{relar}\AgdaSpace{}%
\AgdaBound{n}\AgdaSpace{}%
\AgdaSymbol{→}\AgdaSpace{}%
\AgdaSymbol{(}\AgdaBound{i}\AgdaSpace{}%
\AgdaSymbol{:}\AgdaSpace{}%
\AgdaDatatype{ℕ}\AgdaSymbol{)}\AgdaSpace{}%
\AgdaSymbol{→}\AgdaSpace{}%
\AgdaDatatype{ℕ}\AgdaSpace{}%
\AgdaOperator{\AgdaFunction{\textasciicircum{}}}\AgdaSpace{}%
\AgdaBound{n}\AgdaSpace{}%
\AgdaSymbol{→}\AgdaSpace{}%
\AgdaPrimitive{Prop}\<%
\\
%
\>[2]\AgdaFunction{ℕRel}\AgdaSpace{}%
\AgdaSymbol{\AgdaUnderscore{}}\AgdaSpace{}%
\AgdaInductiveConstructor{eq}\AgdaSpace{}%
\AgdaSymbol{\AgdaUnderscore{}}\AgdaSpace{}%
\AgdaSymbol{(}\AgdaBound{m}\AgdaSpace{}%
\AgdaOperator{\AgdaInductiveConstructor{,}}\AgdaSpace{}%
\AgdaBound{n}\AgdaSpace{}%
\AgdaOperator{\AgdaInductiveConstructor{,}}\AgdaSpace{}%
\AgdaSymbol{\AgdaUnderscore{})}\AgdaSpace{}%
\AgdaSymbol{=}\AgdaSpace{}%
\AgdaBound{m}\AgdaSpace{}%
\AgdaOperator{\AgdaDatatype{≡}}\AgdaSpace{}%
\AgdaBound{n}\<%
\\
%
\\[\AgdaEmptyExtraSkip]%
%
\>[2]\AgdaKeyword{import}\AgdaSpace{}%
\AgdaModule{Family}\AgdaSpace{}%
\AgdaDatatype{funar}\AgdaSpace{}%
\AgdaDatatype{relar}\AgdaSpace{}%
\AgdaDatatype{ℕ}\AgdaSpace{}%
\AgdaSymbol{(λ}\AgdaSpace{}%
\AgdaBound{I}\AgdaSpace{}%
\AgdaSymbol{→}\AgdaSpace{}%
\AgdaDatatype{ℕ}\AgdaSymbol{)}\AgdaSpace{}%
\AgdaFunction{ℕfun}\AgdaSpace{}%
\AgdaFunction{ℕRel}\ as\ \AgdaModule{F}\<%
\\
%
\>[2]\AgdaKeyword{open}\AgdaSpace{}%
\AgdaModule{Model}\AgdaSpace{}%
\AgdaFunction{F.F}\<%
\end{code}
\newcommand{\pFamE}{
\begin{code}%
%
\>[2]\AgdaFunction{A⊃B⊃A}\AgdaSpace{}%
\AgdaSymbol{:}\AgdaSpace{}%
\AgdaSymbol{\{}\AgdaBound{A}\AgdaSpace{}%
\AgdaBound{B}\AgdaSpace{}%
\AgdaSymbol{:}\AgdaSpace{}%
\AgdaFunction{For}\AgdaSpace{}%
\AgdaFunction{◇}\AgdaSymbol{\}}\AgdaSpace{}%
\AgdaSymbol{→}\AgdaSpace{}%
\AgdaFunction{Pf}\AgdaSpace{}%
\AgdaSymbol{(}\AgdaFunction{◇}\AgdaSymbol{)}\AgdaSpace{}%
\AgdaSymbol{(}\AgdaBound{A}\AgdaSpace{}%
\AgdaOperator{\AgdaFunction{⊃}}\AgdaSpace{}%
\AgdaBound{B}\AgdaSpace{}%
\AgdaOperator{\AgdaFunction{⊃}}\AgdaSpace{}%
\AgdaBound{A}\AgdaSymbol{)}\<%
\\
%
\>[2]\AgdaFunction{A⊃B⊃A}\AgdaSpace{}%
\AgdaBound{i}\AgdaSpace{}%
\AgdaBound{γ}\AgdaSpace{}%
\AgdaBound{a}\AgdaSpace{}%
\AgdaBound{b}\AgdaSpace{}%
\AgdaSymbol{=}\AgdaSpace{}%
\AgdaBound{a}\<%
\end{code}}
\newcommand{\pFamC}{
\begin{code}%
%
\>[2]\AgdaFunction{∀P∧Q⊃∀P∧∀Q}\AgdaSpace{}%
\AgdaSymbol{:}%
\>[1500I]\AgdaSymbol{\{}\AgdaBound{P}\AgdaSpace{}%
\AgdaBound{Q}\AgdaSpace{}%
\AgdaSymbol{:}\AgdaSpace{}%
\AgdaFunction{For}\AgdaSpace{}%
\AgdaSymbol{(}\AgdaFunction{◇}\AgdaSpace{}%
\AgdaOperator{\AgdaFunction{▹ₜ}}\AgdaSymbol{)\}}\AgdaSpace{}%
\AgdaSymbol{→}\<%
\\
\>[1500I][@{}l@{\AgdaIndent{0}}]%
\>[16]\AgdaFunction{Pf}\AgdaSpace{}%
\AgdaSymbol{(}\AgdaFunction{◇}\AgdaSymbol{)}\AgdaSpace{}%
\AgdaSymbol{((}\AgdaFunction{Forall}\AgdaSpace{}%
\AgdaSymbol{(}\AgdaBound{P}\AgdaSpace{}%
\AgdaOperator{\AgdaFunction{∧}}\AgdaSpace{}%
\AgdaBound{Q}\AgdaSymbol{))}\AgdaSpace{}%
\AgdaOperator{\AgdaFunction{⊃}}\AgdaSpace{}%
\AgdaSymbol{(}\AgdaFunction{Forall}\AgdaSpace{}%
\AgdaBound{P}\AgdaSpace{}%
\AgdaOperator{\AgdaFunction{∧}}\AgdaSpace{}%
\AgdaFunction{Forall}\AgdaSpace{}%
\AgdaBound{Q}\AgdaSymbol{))}\<%
\\
%
\>[2]\AgdaFunction{∀P∧Q⊃∀P∧∀Q}\AgdaSpace{}%
\AgdaBound{i}\AgdaSpace{}%
\AgdaBound{γ}\AgdaSpace{}%
\AgdaBound{x}\AgdaSpace{}%
\AgdaSymbol{=}\AgdaSpace{}%
\AgdaSymbol{(λ}\AgdaSpace{}%
\AgdaBound{d}\AgdaSpace{}%
\AgdaSymbol{→}\AgdaSpace{}%
\AgdaField{fst}\AgdaSpace{}%
\AgdaSymbol{(}\AgdaBound{x}\AgdaSpace{}%
\AgdaBound{d}\AgdaSymbol{))}\AgdaSpace{}%
\AgdaOperator{\AgdaInductiveConstructor{,p}}\AgdaSpace{}%
\AgdaSymbol{λ}\AgdaSpace{}%
\AgdaBound{d}\AgdaSpace{}%
\AgdaSymbol{→}\AgdaSpace{}%
\AgdaField{snd}\AgdaSpace{}%
\AgdaSymbol{(}\AgdaBound{x}\AgdaSpace{}%
\AgdaBound{d}\AgdaSymbol{)}\<%
\end{code}}

Ebben a fejezetben példákon keresztül szemléltetjük, hogyan lehet konkrét logikai állításokat bizonyítani az egyes megadott modellekben. Egyszerűbb és összetettebb formulák esetén is megvizsgáljuk, hogyan néz ki a bizonyítás a szintaxis szintjén, illetve hogyan értelmezhetők a különböző szemantikai modellekben (Bool, Tarski, Family). Külön hangsúlyt fektetünk arra, hogy az egyes modellekben való bizonyítás mennyiben könnyebb vagy mikben tér el.

\section{Egyszerűbb formula}

Egyszerű formulának a következőt választottuk:
\begin{align*}
\texttt{A ⊃ B ⊃ A}
\end{align*}

\subsection{Informális bizonyítás}

\[
\infer[\text{⊃in}]{\vdash A ⊃ (B ⊃ A)}{
  \infer[\text{⊃in}]{A \vdash B ⊃ A}{
    \infer[\text{\small W-L}]{A, B \vdash A}{
        A \vdash A \quad \text{\small [Ax]}
    }
  }
}
\]

\subsection{Bizonyítás a szintaxisban}

A szintaktikai megközelítés a formális logika alapelveit követi, ahol a bizonyítás lényege a levezetési szabályok alkalmazása. Ez a módszer nem azzal foglalkozik, hogy mit "jelentenek" a logikai állítások, hanem inkább azzal, hogyan lehet őket levezetni más állításokból a rendszer szabályai szerint. Itt láthatjuk a legtöbb hasonlóságot az informális bizonyítással.

\pSynE

\subsection{Bizonyítás a Bool modellben}

A Bool modell a klasszikus igazságérték-alapú megközelítést képviseli. Ebben a modellben a bizonyítás lényege, hogy minden lehetséges igazságérték-kombinációt megvizsgálunk, akárcsak egy igazságtáblán. Agdában tehát e modell esetében sokszor elegendő mintát illeszteni a formulákra.

\pBooE

\subsection{Bizonyítás a Tarski modellben}

A Tarski modellbeli bizonyításnál láthatjuk, hogy az implikáció definíciója miatt - azaz mivel az Agda beépített függvény típusát használtuk - könnyű dolgunk lesz. Bármilyen \textit{γ} kontextusban, ha kapunk egy \textit{a} értéket (amely az A formula γ kontextusbeli értéke), valamint egy \textit{b} értéket (a B formula γ kontextusbeli értéke), akkor egyszerűen visszaadjuk az \textit{a} értéket. Ez pontosan azt fejezi ki, hogy "ha A, akkor (ha B, akkor A)" - vagyis függetlenül B értékétől, az eredmény mindig A értéke lesz. A Tarski modell ereje abban rejlik, hogy matematikai tisztasággal, absztrakt módon ragadja meg a logikai összefüggéseket, ami különösen összetettebb formuláknál válik előnyössé a Bool modellhez képest.

\pTarE

\subsection{Bizonyítás a Family modellben}

A Family modellben megadott bizonyítás nagyon hasonlít a Tarski modellbeli bizonyításhoz, de kapunk egy extra \textit{i} paramétert. A bizonyítás lényegében ugyanazon az elven alapul, mint a Tarski modellben: függetlenül B értékétől, az eredmény mindig A értéke lesz.

\pFamE

\subsection{Bizonyítás bármely modellben}

A "bármely modellben" történő bizonyítás a legáltalánosabb és egyben a legösszetettebb is. Úgy kell megfogalmaznunk a bizonyítást, hogy az minden modellben érvényes legyen. Ez a megközelítés különösen értékes elméleti szempontból, hiszen rávilágít a különböző logikai rendszerek közötti összefüggésekre.

\pAnyE

\section{Összetettebb formula}

Az választott összetettebb formula:
\begin{align*}
\texttt{∀(P ∧ Q) ⊃ ∀P ∧ ∀Q}
\end{align*}

\subsection{Informális bizonyítás}

{\fontsize{10pt}{10pt}\selectfont
\[
\infer[⊃in]{\vdash ∀x(P(x) ∧ Q(x)) ⊃ (∀x\,P(x) ∧ ∀x\,Q(x))}{
  \infer[∧in]{∀x(P(x) ∧ Q(x)) \vdash ∀x\,P(x) ∧ ∀x\,Q(x)}{
    \infer[∀in]{∀x(P(x) ∧ Q(x)) \vdash ∀x\,P(x)}{
      \infer[∧out₁]{∀ x(P(x) ∧ Q(x)) \vdash P(x)}{
        \infer[∀out]{∀ x(P(x) ∧ Q(x)) \vdash P(x) ∧ Q(x)}{
            ∀x(P(x) ∧ Q(x)) \vdash ∀x(P(x) ∧ Q(x) \quad \text{[Ax]}
        }
      }
    }
    &
    \infer[∀in]{∀x(P(x) ∧ Q(x)) \vdash ∀x\,Q(x)}{
      \infer[∧out₂]{∀x(P(x) ∧ Q(x)) \vdash Q(x)}{
        \infer[∀out]{∀x(P(x) ∧ Q(x)) \vdash P(x) ∧ Q(x)}{
            ∀x(P(x) ∧ Q(x)) \vdash ∀x(P(x) ∧ Q(x) \quad \text{[Ax]}
        }
      }
    }
  }
}
\]
}

\subsection{Bizonyítás a szintaxisban}

Itt szintén láthatjuk, hogy a szintaxisbeli bizonyítás tükrözi a fenti informális levezetést.

\pSynC

\subsection{Bizonyítás a Bool modellben}

Ennek a megközelítésnek az előnye, hogy pontosan látjuk, hogyan működik a formula minden esetben. A hátránya viszont nyilvánvalóvá válik, amikor összetettebb formulákkal dolgozunk. Jelen példa is rendkívül hosszú és bonyolult, tele van esetszétválasztásokkal és feltételes állításokkal. Ez jól mutatja, hogy bár a Bool modell intuitív, a gyakorlatban nehézkessé válhat.

\pBooC

\subsection{Bizonyítás a Tarski modellben}

A Tarski modell tiszta matematikai absztrakciója itt különösen jól látható. A bizonyítás a metaelmélet eszközeit használja, ami jelentősen egyszerűsíti a munkát. Lehetővé teszi, hogy ahelyett, hogy minden eseten végigmennénk (mint a Bool modellben), egyszerűen megadjuk azt a függvényt, amely az állítás igazságát konstruálja. A metaelmélet konstruktorait alkalmazva közvetlenül építhetjük fel a bizonyítást, ami tömörebb és átláthatóbb eredményt ad.

\pTarC

\subsection{Bizonyítás a Family modellben}

A Family modellben az összetettebb formula bizonyítása szintén nagyon hasonlít a Tarski modellbeli megfelelőjére, de itt is megjelenik az extra \textit{i} paraméter.

\pFamC

\subsection{Bizonyítás bármely modellben}

Az összetettebb formulánál az \textit{any model} bizonyításunk nagyon absztrakt és számos transformációt (substP) használ, mert arra törekszik, hogy minden lehetséges szemantikai értelmezésben egyaránt érvényes legyen a bizonyítás. A kód komplexitása jól mutatja, milyen kihívásokkal jár egy teljesen általános bizonyítás megalkotása, amely nem támaszkodhat egyik specifikus modell előnyeire sem. Ugyanakkor ez a megközelítés biztosítja, hogy a bizonyítás valóban univerzális legyen.

\pAnyC


\section{Összehasonlítás}

A bemutatott példák jól szemléltetik, hogy a különböző szemantikák más-más filozófiai alapokon nyugszanak, de ugyanarra az eredményre vezetnek, valamint a bizonyítások formája és komplexitása eltér. A szintaktikai bizonyítás a formális szabályokat hangsúlyozza, a Bool modell az igazságértékekre koncentrál, míg a Tarski és Family modellek absztraktabb, matematikai struktúrákat követnek.

A bármely modellben történő bizonyítás a legjobb olyan szempontból, hogy ha itt sikerül megadnunk egy bizonyítást, akkor az minden lehetséges modellben érvényes lesz. Ez az univerzalitás azonban komoly nehézségekkel jár a gyakorlati kivitelezés során. A fő probléma abban rejlik, hogy az egyenlőségek ebben a megközelítésben pusztán posztulált formában jelennek meg, nem pedig definíció szerint teljesülnek. Ez jelentősen megnehezíti a bizonyítások kidolgozását.

A szintaxisban történő bizonyítás már kedvezőbb, mivel itt már legalább a helyettesítési egyenlőségek definíció szerint teljesülnek. Ugyanakkor korlátos abban az értelemben, hogy még nem adtuk meg az inicialitását, így jelenleg nem elég praktikus a használata összetettebb esetekben. Az inicialitás kidolgozása további kutatási feladat marad.

A Tarski modellel a legkönnyebb dolgoznunk, mivel itt minden egyenlőség definíció szerint teljesül, csak a metaelmélet eszköztárát tudjuk használni. Hátránya viszont, hogy a Tarski modell nem biztosít automatikus konverziót más modellekbe. Ezért ha egy formulát bizonyítottunk a Tarski modellben, az nem jelenti azt, hogy rendelkezünk a formula bizonyításával például a Bool modellben is.